\documentclass[11pt,a4paper]{article}
\usepackage[T1]{fontenc}
\usepackage[utf8]{inputenc}
\usepackage[margin=25mm]{geometry}
\usepackage{amsmath,amssymb,amsthm,mathtools}
\usepackage{array,longtable,booktabs,enumitem}
\usepackage[hidelinks]{hyperref}
\usepackage{xurl}
\setlength{\emergencystretch}{3em}
\setlength{\parskip}{4pt}
\newtheorem{theorem}{Theorem}[section]
\newtheorem{lemma}[theorem]{Lemma}
\newtheorem{proposition}[theorem]{Proposition}
\newtheorem{corollary}[theorem]{Corollary}
\theoremstyle{definition}
\newtheorem{definition}[theorem]{Definition}
\newtheorem{remark}[theorem]{Remark}
\newtheorem{example}[theorem]{Example}
\newcommand{\R}{\mathbb R}
\newcommand{\C}{\mathbb C}
\newcommand{\N}{\mathbb N}
\newcommand{\F}{\mathbb F}
\newcommand{\E}{\mathbb E}
\newcommand{\norm}[1]{\left\lVert#1\right\rVert}
\newcommand{\abs}[1]{\left|#1\right|}
\newcommand{\inner}[2]{\left\langle#1,#2\right\rangle}
\newcommand{\op}{\mathrm{op}}
\DeclareMathOperator{\diag}{diag}
\DeclareMathOperator{\tr}{tr}
\DeclareMathOperator{\rank}{rank}
\DeclareMathOperator{\supp}{supp}
\DeclareMathOperator{\ran}{ran}
\DeclareMathOperator{\Span}{span}
\DeclareMathOperator{\dist}{dist}
\DeclareMathOperator{\conv}{conv}
\newcommand{\code}[1]{\ifmmode\text{\small\nolinkurl{#1}}\else{\small\nolinkurl{#1}}\fi}
\title{Remaining mathematical items in NLAlib\\
Finite linear algebra, operator proofs, and a statement audit}
\author{Mathematical handoff for subsequent Lean formalization}
\date{13 October 2026}
\begin{document}
\maketitle

\section*{Scope and evidence}
The refreshed public site and the repository agree at commit
\href{https://github.com/tchen-research/NLAlib/tree/e5725bb94e5bd8e326d58ad219260b0218e74c9e}
{\texttt{e5725bb94e5b}}: 291 entries are marked proved and 63 are open
(56 candidate, 6 scaffold, 1 assumed). All 108 matrix-analysis-toolkit
entries are marked proved. This supersedes the earlier list of twelve
matrix candidates; those were a catalogue-status mismatch.

This note selects \textbf{33 mathematical entries} across the catalogue.
It includes structural algebra, probability identities, approximation
theory, and quadrature. The 30 excluded entries concern algorithm
definitions, recurrences, implementation identities, query models,
or algorithm-specific convergence and performance guarantees; they
are listed at the end so the scope decision is reviewable.
The unbiasedness entries are included for their expectation identities,
and interpolative decomposition for its existence theorem.

The purpose is to give mathematics that another agent can formalize.
No Lean source is changed, no axiom is proposed, and no new kernel
verification is claimed. Declaration names and signatures were checked
against the commit-pinned \code{atlas/declarations.json}; catalogue
contracts were read from \code{atlas/atlas.json}. Additional Mathlib
names were inspected in the local checkout at commit
\texttt{0df444a360ea}. These checks identify reusable interfaces;
they are not a fresh build of the remote repository.

\paragraph{Read the status of each route literally.}
Most sections give complete mathematical reductions to standard finite
algebra and scalar analysis. The Crouzeix--Palencia section explicitly
retains its scalar Cauchy-transform boundary theorem. The two sharp
approximation entries still retain the Stahl and Halphen strong
asymptotics; finite certificates are not presented as proofs of those
asymptotics. The unspecified early-truncation identity cannot be
reconstructed from its catalogue title. Several other advertised
clauses need hypothesis or sign repairs, which are stated where used.

\paragraph{Shared conventions.}
Matrix norms without a subscript are Euclidean operator norms, and
$\|\cdot\|_F$ is Frobenius norm. Adjoint means conjugate transpose;
in real statements it is transpose. Empty sums and products are $0$
and $1$, respectively. Measures in positive orthogonality statements
are nonzero finite positive Borel measures with the required finite
moments. A finite target of order $n$ requires only finitely many
moments, not an infinite-dimensional polynomial completion.
For monic orthogonal polynomials, $h_j=\int\pi_j^2\,d\mu$,
$\beta_j=h_j/h_{j-1}$, and the symmetric Jacobi off-diagonal is
$a_j=\sqrt{\beta_j}$; normalization is made explicit in each section.

\tableofcontents
\clearpage


\section{Positive real moments: normalization and one scalar tangent}
\label{math:moment}
\paragraph{Catalogue item.} \code{positive-real-moment-comparison}.
Let $(\Omega,\mathcal F,\mu)$ be any probability space, let $f\ge0$
be almost everywhere strongly measurable, and suppose
$0<p\le q$ and $f^q$ is integrable. Set $M=\int f^q\,d\mu$. Then
\[
 \left(\int f^p\,d\mu\right)^{1/p}\le M^{1/q}.
\]
No assumption $p\ge1$ is needed.

\begin{proof}
The case $p=q$ is equality. Otherwise set $a=p/q\in(0,1)$.
The scalar tangent inequality
\[
 t^a\le (1-a)+at\qquad(t\ge0)
\]
follows by differentiating on $(0,\infty)$ and checking $t=0$.
If $M=0$, positivity gives $f^q=0$ almost everywhere, hence $f=0$
almost everywhere. If $M>0$, apply the inequality to $g=f^q/M$.
It first proves integrability of $g^a$, by domination by $1-a+ag$,
and then gives $\int g^a\le1$. Thus $\int f^p\le M^{p/q}$;
raising to the positive power $1/p$ proves the assertion.
\end{proof}

\paragraph{Actual reuse.}
Mathlib already contains
\code{MeasureTheory.eLpNorm'_le_eLpNorm'_of_exponent_le},
\code{MeasureTheory.eLpNorm_le_eLpNorm_of_exponent_le}, and
\code{MeasureTheory.MemLp.mono_exponent} in
\code{MeasureTheory/Function/LpSeminorm/CompareExp.lean}.
The first takes real exponents with $0<p\le q$, including exponents
below one. The rooted integral formulation is a conversion using
\code{MeasureTheory.MemLp.eLpNorm_eq_integral_rpow_norm}.
The direct proof above is useful if avoiding extended-nonnegative-real
conversions makes the wrapper shorter; it is not a missing probability theorem.

\section{Random projections and unbiased trace identities}
\label{math:trace}
\paragraph{Catalogue items.}
\code{hutchpp-unbiased}, \code{xtrace-unbiased}.
These are expectation identities; no estimator construction or
iteration analysis is required.

Let $A\in\R^{d\times d}$ be arbitrary, including nonsymmetric $A$.
Let $Q^\top Q=I$, put $P=QQ^\top$, and write
\[
 D(Q,w)=\tr(Q^\top AQ)+w^\top(I-P)A(I-P)w.
\]
For an isotropic probe $w$ we assume genuine finite second moments and
$\mathbb E(w_iw_j)=\delta_{ij}$. In particular
$\mathbb E\|w\|^2=d$. Finite summation gives
\[
 \mathbb E(w^\top Bw)=\tr B.
\]
The scalar terms are integrable because
$2|w_iw_j|\le w_i^2+w_j^2$. Cyclicity and $P^2=P$ give
\[
 \tr(Q^\top AQ)+\tr((I-P)A(I-P))
 =\tr(AP)+\tr(A(I-P))=\tr A.                 \tag{T1}
\]

\begin{lemma}[Uniform domination removes the integrability side condition]
For every orthonormal $Q$ and every $w$,
\[
 |D(Q,w)|\le \|A\|_2\bigl(d+\|w\|^2\bigr). \tag{T2}
\]
\end{lemma}
\begin{proof}
The trace term is the sum of at most $d$ numbers $q_j^\top Aq_j$,
each bounded in absolute value by $\|A\|_2$. Orthogonal projection
is contractive, so the other term is at most $\|A\|_2\|w\|^2$.
The proof also covers empty columns and $d=0$.
\end{proof}

Suppose $S$ is any random element, $Q(S)$ is measurable and has
orthonormal columns almost surely, and $S$ and $w$ are independent.
Their joint law is the product of their marginal laws. Bound (T2)
proves integrability on that product law. Integrate first in $w$,
use (T1), and integrate in $S$:
\[
 \mathbb E D(Q(S),w)=\tr A.                  \tag{T3}
\]
This uses neither a conditional distribution nor a standard Borel
assumption on the underlying probability space. A finite nonempty
average of such terms has the same mean; mutual independence between
the terms is unnecessary.

For leave-one-out probes $w_1,\ldots,w_m$, assume the family is
independent and each probe is measurable and isotropic. Let
$S_i=(w_j)_{j\ne i}$ and $Q_i=F_i(S_i)$ for a measurable $F_i$.
Independence of disjoint subfamilies gives $S_i\perp w_i$;
composition preserves independence, and (T2)--(T3) apply to each $i$.
Identical distribution is sufficient but unnecessary. Exchangeability
alone would not establish this independence.

\paragraph{Actual reuse and the remaining wrapper.}
The snapshot already has
\code{NLAlib.integral_deflatedTraceEstimate_eq_trace},
\code{NLAlib.integral_deflatedTraceEstimate_comp_eq_trace}, and
\code{NLAlib.integral_leaveOneOutDeflatedTraceEstimate_eq_trace}.
The last two require product-law integrability explicitly; (T2)
discharges it for a measurable orthonormal correction.
Mathlib's \code{ProbabilityTheory.iIndepFun.indepFun_finset}
(and its almost-everywhere version)
supplies held-out independence. This is the precise missing iid
construction bridge recorded in the refreshed catalogue.

\section{Interpolative decomposition without strong RRQR}
\label{math:id}
\paragraph{Catalogue item.} \code{interpolative-decomposition}.
Let $A\in\R^{m\times n}$ and $0\le k\le\rank A$.
There exist $k$ distinct column indices $J$ and
$T\in\R^{k\times n}$ such that
\[
 T[:,J]=I_k,\qquad |T_{ij}|\le1,\qquad
 \|A-A[:,J]T\|_2\le\sqrt{1+k(n-k)}\,\sigma_{k+1}(A). \tag{ID}
\]
Here singular values past the smaller dimension are zero.
If $\rank A=k$, this is an exact factorization.

\begin{proof}
For $k=0$ take empty $J,T$; (ID) is $\|A\|_2=\sigma_1(A)$.
For $k>0$, take an SVD and let $V\in\R^{k\times n}$ be its first
$k$ right singular rows, so $VV^\top=I_k$.
Choose $J$ maximizing $|\det V[:,J]|$ among the finitely many
$k$-subsets. Some minor is nonzero because $V$ has row rank $k$.
Define
\[
 T=V[:,J]^{-1}V.
\]
Cramer's rule expresses $T_{ij}$ as the ratio of a one-column
replacement determinant to the maximal determinant. Replacement by
an already selected column either gives zero or the same minor;
otherwise it gives a $k$-subset minor up to sign. Hence $|T_{ij}|\le1$.

Let $A_k$ be the rank-$k$ SVD truncation, $F=A-A_k$, and let $S_J$
be the coordinate inclusion selecting columns $J$. Since $A_k=BV$
for an $m\times k$ matrix $B$,
\[
 A_kS_JT=A_k,\qquad
 A-AS_JT=F(I-S_JT).                          \tag{ID1}
\]
After moving $J$ first, write $T=[I_k\ T_2]$. Then
\[
 I-S_JT=\begin{pmatrix}0&-T_2\\0&I_{n-k}\end{pmatrix},
 \quad
 \|I-S_JT\|_2\le\sqrt{1+\|T_2\|_2^2}
 \le\sqrt{1+k(n-k)}.
\]
The first inequality follows immediately by applying the block
matrix to $(x,y)$; no norm formula at dimension zero is needed.
Combine this with $\|F\|_2=\sigma_{k+1}(A)$ in (ID1).
If $k=n$, $T_2$ has no columns and the residual is exactly zero.
\end{proof}

\paragraph{Why this route is shorter.}
Maximize a minor of the \emph{right singular frame}, not the volume
of the original columns. No singular-value estimates for a QR trailing
block are needed. The selected columns of $A$ are independent because
their images under the top left singular vectors have invertible
matrix $\diag(\sigma_1,\ldots,\sigma_k)V[:,J]$.
This is an existence proof; the catalogue's additional claim about an
efficient coefficient-$2$ construction belongs to algorithmic scope.

\paragraph{Actual reuse.}
\code{NLAlib.exists_svd_right_blocks_specNorm} already supplies an
orthogonal right-frame decomposition and the exact tail norm.
The additional lemma is only finite maximal-minor interpolation.
The same proof gives a Frobenius bound with
$\|A-A_k\|_F$ in place of $\sigma_{k+1}(A)$.

\section{Projected truncation: a complete identity and a precise scope limit}
\label{math:truncation}
\paragraph{Catalogue item.} \code{truncation-lemma}.
For every $Q$ with orthonormal columns and every compatible $X$,
orthogonality gives the exact identity
\[
 \|A-QX\|_F^2
 =\|(I-QQ^\top)A\|_F^2+\|Q^\top A-X\|_F^2. \tag{PT}
\]
Indeed the two summands in
$A-QX=(I-QQ^\top)A+Q(Q^\top A-X)$ have zero Frobenius inner
product, and multiplication by $Q$ preserves the Frobenius norm.
Take $X$ to be any best rank-at-most-$k$ approximation of $Q^\top A$.
If $A_k$ is a best rank-at-most-$k$ approximation of $A$, then
$Q^\top A_k$ is an admissible competitor, so
\[
 \|A-QX\|_F^2\le\|(I-QQ^\top)A\|_F^2+\|A-A_k\|_F^2.
\]
This is stronger than the advertised coefficient $2$ on the projection
residual, and includes $k=0$ and empty matrices.

\paragraph{Do not invent the unspecified component.}
This ordinary statement is already proved under
\code{truncation-lemma-projected}, including
\code{NLAlib.IsBestFrobRankApprox.frobSq_sub_mul_le_bestRankFrobSq_add_frobSq_residual}.
The remaining catalogue text mentions an ``early-truncated variant''
without defining its output or giving its identity. Its two historical
external declaration pointers are explicitly unverified in the atlas.
Equation (PT) is available for any specified $X$, but does not identify
an unspecified algorithm output. That clause is not counted as closed
by the present derivation.


\section{Vector-relative minimal polynomials and Krylov grade}
\label{struct:grade}
\paragraph{Catalogue target.} \code{krylov-grade}. The remaining extension is
over an arbitrary field, with a vector-relative annihilator and a divisibility
bridge. No inner product, spectral theorem, or principal-ideal construction is
needed for the following route.

Let $V$ be an $n$-dimensional vector space over a field $\F$, let $A:V\to V$
be linear, and put
\[
 K_q(A,b)=\Span\{b,Ab,\ldots,A^{q-1}b\},\qquad K_0=\{0\}.
\]
Choose the least integer $\nu\ge0$ for which $A^\nu b\in K_\nu$.
Such an integer exists with $\nu\le n$: otherwise the first $n+1$ powers
are independent, by taking the highest nonzero coefficient of a relation.
For every $q\le\nu$, that same argument proves independence of
$b,\ldots,A^{q-1}b$. Thus there are unique $c_0,\ldots,c_{\nu-1}$ such that
\[
 A^\nu b=-\sum_{j<\nu}c_j A^j b.
\]
Define $\mu_b(t)=t^\nu+\sum_{j<\nu}c_jt^j$; when $\nu=0$, this is $1$
and $b=0$.

\begin{lemma}[Division is the whole annihilator interface]
For every $p\in\F[t]$,
\[
 p(A)b=0\quad\Longleftrightarrow\quad\mu_b\mid p.
 \tag{S1}
\]
In particular, $\mu_b$ is the unique monic polynomial of least degree
annihilating $b$, and generates the vector-annihilator ideal.
\end{lemma}
\begin{proof}
Divide $p=q\mu_b+r$ with $\deg r<\nu$. Polynomial evaluation commutes with
$A$, so $p(A)b=r(A)b$. Independence of the first $\nu$ powers forces all
coefficients of $r$ to vanish. Conversely a multiple of $\mu_b$ annihilates
$b$. For $\nu=0$, division by $1$ has zero remainder and every polynomial
annihilates the zero vector.
\end{proof}

\begin{theorem}[All grade conclusions over a field]
For every $q\ge0$,
\[
 \dim K_q=\min(q,\nu),\qquad
 (b,Ab,\ldots,A^{q-1}b)\text{ independent}\ \Longleftrightarrow\ q\le\nu.
\]
Moreover $AK_\nu\subseteq K_\nu$, $K_q=K_\nu$ for $q\ge\nu$,
$\mu_b\mid\mu_A$, and $\nu\le\deg\mu_A\le n$.
If $A$ is invertible, $A^{-1}b\in K_\nu$.
\end{theorem}
\begin{proof}
Applying $A$ to each spanning vector of $K_\nu$ remains in that space:
only its last vector uses the defining relation. Induction therefore puts
every later power in $K_\nu$, proving saturation and the dimension formula.
The matrix minimal polynomial $\mu_A$ annihilates $b$, so (S1) gives
divisibility. Cayley--Hamilton gives $\mu_A\mid p_A$ and hence the degree
bound. If $b\ne0$ and $A$ is invertible, the constant coefficient
$c_0$ is nonzero. Otherwise $\mu_b=tq$ and
$0=Aq(A)b$ would imply $q(A)b=0$, contradicting minimal degree.
Multiplying the defining relation by $A^{-1}$ gives the explicit witness
\[
 A^{-1}b=-c_0^{-1}\bigl(A^{\nu-1}b+c_{\nu-1}A^{\nu-2}b+\cdots+c_1b\bigr).
\]
For $b=0$ membership is immediate. Dimension zero is included.
\end{proof}

\paragraph{Verified reuse and the actual new interface.}
The snapshot already has \code{NLAlib.mem_krylovSpace_iff} and
\code{NLAlib.finrank_krylovSpace_le} over an arbitrary field. Thus the
catalogue note that \code{krylovSpace} itself still needs generalization
is stale. The indexed \code{NLAlib.krylovGrade},
\code{NLAlib.finrank_krylovSpace},
\code{NLAlib.linearIndependent_pow_mulVec_iff_le_krylovGrade}, and
\code{NLAlib.krylovSpace_eq_of_krylovGrade_le} are still real statements.
The reusable addition is (S1), followed by the field-general grade bridge.
This avoids constructing a new polynomial-module annihilator merely to
recover polynomial division.

\section{Companion matrices and cyclic vectors without rational canonical form}
\label{struct:companion}
\paragraph{Catalogue target and correction.} \code{companion-matrix}.
For a real matrix, ``every eigenvalue geometrically simple'' must refer
to eigenvalues and eigenspaces over $\C$, or more generally a splitting
field. Real eigenvalues alone do not suffice: two copies of
$\left(\begin{smallmatrix}0&-1\\1&0\end{smallmatrix}\right)$ have no real
eigenvalue, but their minimal polynomial has degree $2<4$.
The algebraic equivalences below hold over every field; the geometric
criterion is stated over a splitting field.

For $p(t)=t^n+c_{n-1}t^{n-1}+\cdots+c_0$, define $C_p$ by
\[
 C_pe_j=e_{j+1}\ (0\le j<n-1),\qquad
 C_pe_{n-1}=-\sum_{j<n}c_je_j.
 \tag{S2}
\]
For $n=0$, $p=1$ and $C_p$ is the empty matrix.
The first $n$ Krylov vectors of $e_0$ are the standard basis when $n>0$.
Equation (S2) gives $p(C_p)e_0=0$, hence $p(C_p)C_p^je_0=0$ for every
basis vector. No polynomial of degree below $n$ annihilates $e_0$.
Thus $\mu_{C_p}=p$. Cayley--Hamilton and equality of monic degrees give
$p_{C_p}=p$ without a determinant expansion.

\begin{lemma}[A vector realizes the operator minimal polynomial]
Every finite-dimensional operator over any field has a vector $b$ with
$\mu_b=\mu_A$.
\end{lemma}
\begin{proof}
The zero-dimensional case uses $b=0$ and $\mu_A=1$. Otherwise factor
$\mu_A=\prod_{j=1}^s f_j^{e_j}$ into distinct monic irreducibles.
Since $\mu_A/f_j$ has smaller degree, $(\mu_A/f_j)(A)$ is not the zero
operator. Choose $w_j$ on which it is nonzero, and set
\[
 b_j=(\mu_A/f_j^{e_j})(A)w_j.
\]
Then $f_j^{e_j}(A)b_j=0$ but $f_j^{e_j-1}(A)b_j\ne0$.
By (S1), $\mu_{b_j}=f_j^{e_j}$. Put $b=\sum_jb_j$.
If $p(A)b=0$, apply $(\mu_A/f_j^{e_j})(A)$: it kills every $b_l$ with
$l\ne j$, and yields
$f_j^{e_j}\mid p\,(\mu_A/f_j^{e_j})$.
Coprimality implies $f_j^{e_j}\mid p$ for all $j$, hence $\mu_A\mid p$.
The converse holds because $\mu_A(A)=0$. Apply (S1) once more.
\end{proof}

\begin{theorem}[The companion and centralizer characterizations]
For $A\in\F^{n\times n}$, the following are equivalent:
\[
 \deg\mu_A=n;\qquad \mu_A=p_A;\qquad
 A\text{ is similar to }C_{p_A};\qquad
 \text{some }b\text{ has grade }n.
 \tag{S3}
\]
When these hold, every $B$ commuting with $A$ is $q(A)$ for a unique
polynomial $q$ of degree $<n$ (with $q=0$ in dimension zero).
\end{theorem}
\begin{proof}
The first two statements are equivalent because $\mu_A\mid p_A$ and
both are monic. The preceding lemma supplies a grade-$n$ vector if
$\deg\mu_A=n$. A grade-$n$ vector is a basis generator; in the basis
$(b,Ab,\ldots,A^{n-1}b)$, its annihilating relation is exactly (S2).
Conversely the companion basis is cyclic and similarity preserves
minimal polynomials. If $AB=BA$, expand $Bb=q(A)b$ in that Krylov basis.
Then $BA^jb=A^jBb=q(A)A^jb$ for all basis vectors, so $B=q(A)$.
Uniqueness follows by applying to $b$ and using independence.
\end{proof}

\paragraph{The geometric criterion, with its finite proof.}
Pass to a field where $p_A$ splits. The degree of the minimal polynomial
is unchanged by scalar extension: the largest independent initial
segment of $I,A,A^2,\ldots$ is unchanged, since the rank of their
coordinate matrix is determined by nonzero minors.
Write the primary spaces as $V_\lambda=\ker(A-\lambda I)^n$.
Bezout identities between distinct powers of $t-\lambda$ give
$V=\bigoplus_\lambda V_\lambda$; on $V_\lambda$ write
$A=\lambda I+N_\lambda$, with $N_\lambda$ nilpotent.
This decomposition requires only polynomial Bezout: for coprime $f,g$
and $(fg)(A)=0$, an identity $uf+vg=1$ writes each vector as a sum
from $\ker f(A)$ and $\ker g(A)$, with zero intersection; iterate.

Let $m=\dim V_\lambda$ and let $e$ be the nilpotency index of $N=N_\lambda$.
Its characteristic polynomial is $t^m$: choose a basis adapted to the
increasing kernels of its powers, so its matrix is strictly triangular.
If $\dim\ker N=1$, the map
\[
 \ker N^{j+1}/\ker N^j\longrightarrow\ker N,\qquad [x]\longmapsto N^jx
\]
is injective, so successive dimensions increase by at most one.
They cannot stop increasing before reaching the whole space: equality
$\ker N^j=\ker N^{j+1}$ propagates to all later powers.
Consequently $e=m$. Conversely, if $e=m$, choose $v$ with $N^{m-1}v\ne0$.
The vectors $v,Nv,\ldots,N^{m-1}v$ are independent (apply a suitable
power of $N$ to the first nonzero coefficient of a relation), hence
form a basis, and $\ker N$ is one-dimensional.
Now
\[
 \deg\mu_A=\sum_\lambda e_\lambda\le
 \sum_\lambda\dim V_\lambda=n,
\]
with equality exactly when each eigenspace has dimension one.
This proves the repaired geometric equivalence in (S3), including
repeated eigenvalues. It avoids Jordan form and the module structure theorem.

\paragraph{Verified reuse.}
The field-general \code{NLAlib.mem_krylovSpace_iff} supplies the cyclic
span interface. The local Mathlib file
\code{Mathlib/RingTheory/Polynomial/Basic.lean} contains
\code{Polynomial.sup_ker_aeval_eq_ker_aeval_mul_of_coprime}, matching
the Bezout decomposition above. No companion declaration appears among
the target's formalizations; (S2), (S1), and the cyclic-basis identity
are the small interfaces to add, rather than a full canonical-form API.

\section{Residual-polynomial minimax: quotient out invisible coefficients}
\label{struct:residual}
\paragraph{Catalogue target.} \code{residual-polynomial}.
Let $\mathcal P_k^0=\{p\in\R[t]:\deg p\le k,\ p(0)=1\}$.
For nonempty compact $S\subset\C$, set
\[
 M_S(p)=\max_{z\in S}|p(z)|,\qquad
 \varepsilon_k(S)=\inf_{p\in\mathcal P_k^0}M_S(p).
\]
Set $M_\varnothing(p)=\varepsilon_k(\varnothing)=0$.
For arbitrary sets, use the supremum in $[0,\infty]$; the constant
polynomial $1$ ensures $0\le\varepsilon_k(S)\le1$.

\paragraph{Residual algebra and the invertibility distinction.}
Put $r_0=b-Ax_0$. Writing $x=x_0+q(A)r_0$ with $\deg q<k$ gives
\[
 b-Ax=(I-Aq(A))r_0=p(A)r_0,\qquad p(t)=1-tq(t).
 \tag{S4}
\]
Conversely $p(0)=1$ implies $p=1-tq$ with $\deg q<k$; this constructs
an affine-Krylov representative for every residual polynomial.
For a \emph{specified} $x$, equality of residuals with that representative
only implies their difference lies in $\ker A$. Thus invertibility is
needed to recover membership of that particular $x$; $A=0$ shows why.
If $Ax_\star=b$ and $x=x_0+q(A)r_0$, however,
\[
 x-x_\star=p(A)(x_0-x_\star)
 \tag{S5}
\]
holds even for singular $A$, because $r_0=-A(x_0-x_\star)$.
The catalogue's suggestion that (S5) itself requires invertibility is
unnecessary. The identity uses the polynomial constructed in (S4);
an arbitrary residual identity alone still cannot determine the error
when $A$ is singular. For $k=0$, take $q=0$ and $p=1$.

\begin{theorem}[Monotonicity and attainment]
For all sets for which the extended supremum convention is used,
\[
 \varepsilon_{k+1}(S)\le\varepsilon_k(S),\qquad
 S\subseteq T\Longrightarrow\varepsilon_k(S)\le\varepsilon_k(T).
\]
If $S$ is compact, including finite or empty $S$, the infimum is attained.
\end{theorem}
\begin{proof}
The two inequalities follow respectively by enlarging the polynomial
class and by increasing each polynomial's objective. For attainment,
the empty case is immediate. Write
$p_a(z)=1+\sum_{j=1}^k a_jz^j$ with $a\in\R^k$, and let
\[
 N=\{a:\ \sum_{j=1}^k a_jz^j=0\text{ for every }z\in S\}.
\]
Choose the Euclidean orthogonal complement $N^\perp$.
Projecting $a$ onto this complement does not change $p_a$ on $S$.
There are finitely many $z_1,\ldots,z_m\in S$ such that
\[
 N=\{a:\ \sum_j a_jz_l^j=0\text{ for }1\le l\le m\}.
 \tag{S6}
\]
Indeed, start with $\R^k$; whenever the current intersection exceeds
$N$, choose a vector outside $N$ and a point detecting it. Intersecting
with that point's evaluation kernel strictly lowers the dimension.
After at most $k$ steps the intersection is $N$.
If $N^\perp=\{0\}$, $p=1$ already minimizes. Otherwise (S6) and
compactness of its unit sphere give $c>0$ with
\[
 c\|a\|\le\max_l\left|\sum_j a_jz_l^j\right|
 \quad(a\in N^\perp).
 \tag{S7}
\]
The minimum is strictly positive on that sphere because evaluation
there is injective. On the nonempty sublevel $M_S(p_a)\le1$,
(S7) gives $\|a\|\le2/c$.
Choose $R\ge\max_{z\in S}|z|$. The explicit estimate
\[
 |M_S(p_a)-M_S(p_{a'})|
 \le\sum_{j=1}^k|a_j-a'_j|R^j
 \tag{S8}
\]
proves continuity. Therefore the sublevel inside $N^\perp$ is closed
and bounded in a finite-dimensional real space, hence compact.
It contains $a=0$, and a continuous function attains a minimum there.
Outside it the objective exceeds $1$, so this is a global minimum.
\end{proof}

The point of (S6)--(S8) is that the original coefficient sublevel can
be unbounded when $S$ has few points. Compactness must be applied
after removing coefficients invisible on $S$, rather than asserted
for all coefficients.

\paragraph{The sharp interval value.}
For $0<a<b$, put $c=(b+a)/(b-a)>1$. Then
\[
 \varepsilon_k([a,b])=\frac1{T_k(c)},\qquad
 p_*(x)=\frac{T_k((a+b-2x)/(b-a))}{T_k(c)}.
 \tag{S9}
\]
For $k=0$, both statements read $1=1$.
For $k\ge1$, $p_*(0)=1$ and its values at
$x_j=(a+b-(b-a)\cos(j\pi/k))/2$, $0\le j\le k$, are
$(-1)^j/T_k(c)$. A competitor of smaller sup norm makes
$p-p_*$ have alternating strict signs at these $k+1$ points,
so it has $k$ distinct roots between them. It also vanishes at
$0<a$, yielding $k+1$ distinct roots despite degree at most $k$.
This is impossible; the upper bound follows from $|T_k|\le1$ on
$[-1,1]$.

\paragraph{Verified reuse.}
The indexed
\code{NLAlib.sub_mem_krylovSpace_iff_exists_sub_mulVec_eq} has the
required \code{IsUnit A} for the specified-iterate converse.
\code{NLAlib.exists_sub_eq_aeval_mulVec_of_sub_mem_krylovSpace}
already proves (S5) without invertibility, and
\code{NLAlib.isLeast_chebyshev_minimax} already supplies (S9).
The remaining minimax interface is the finite coefficient argument
(S6)--(S8) and elementary monotonicity.

\section{Numerical range: a sphere projected onto a disk}
\label{struct:fov}
\paragraph{Catalogue target.} \code{field-of-values-def}.
For a complex matrix $A$ on an $n$-dimensional Euclidean space, define
\[
 W(A)=\{x^*Ax:\ \|x\|=1\}.
\]
For real matrices this means the numerical range of the complexification.
If $n=0$, $W(A)=\varnothing$. All set statements below still hold then;
the interval using smallest and largest eigenvalues explicitly requires
$n>0$.

\begin{lemma}[The two-dimensional compression lemma]
The numerical range of every complex $2\times2$ matrix is convex.
\end{lemma}
\begin{proof}
For a unit vector $x=(u,v)$ set
\[
 r_1=2\operatorname{Re}(\bar u v),\quad
 r_2=2\operatorname{Im}(\bar u v),\quad r_3=|u|^2-|v|^2.
\]
Direct expansion gives $r_1^2+r_2^2+r_3^2=1$ and
\[
 x^*Bx=\tfrac12\tr B+\tfrac12(b_{12}+b_{21})r_1
       +\tfrac{i}{2}(b_{12}-b_{21})r_2
       +\tfrac12(b_{11}-b_{22})r_3.
 \tag{S10}
\]
Conversely every $r\in S^2\subset\R^3$ comes from such a vector:
if $r_3>-1$, use
$u=\sqrt{(1+r_3)/2}$ and $v=(r_1+ir_2)/(2u)$;
if $r_3=-1$, use $(u,v)=(0,1)$.
Thus $W(B)=\tfrac12\tr B+L(S^2)$ for a real-linear map
$L:\R^3\to\C\simeq\R^2$.

Rank--nullity gives a unit $h\in\ker L$. If $\|r\|\le1$, choose
\[
 t=-\langle r,h\rangle+
       \sqrt{\langle r,h\rangle^2+1-\|r\|^2}.
\]
Then $\|r+th\|=1$ and $L(r+th)=L(r)$. Hence
$L(S^2)=L(\overline B(0,1))$. A linear image of a convex ball is
convex. This proof handles rank-zero and rank-one maps as well as
genuine ellipses, without classifying ellipse axes.
\end{proof}

\begin{theorem}[All numerical-range properties]
$W(A)$ is compact and convex, contains every eigenvalue, and is
contained in $\{z:|z|\le\|A\|\}$. Moreover
\begin{align*}
 W(\alpha A+\beta I)&=\alpha W(A)+\beta,\\
 W(U^*AU)&=W(A)\quad(U\text{ unitary}),\\
 W(Q^*AQ)&\subseteq W(A)\quad(Q^*Q=I).
\end{align*}
For $n>0$ and $H=(A+A^*)/2$,
\[
 \operatorname{Re}W(A)=[\lambda_{\min}(H),\lambda_{\max}(H)].
 \tag{S11}
\]
If $A$ is normal, $W(A)=\conv\sigma(A)$.
\end{theorem}
\begin{proof}
Compactness is the continuous image of the finite-dimensional unit
sphere. A normalized eigenvector gives its eigenvalue in $W(A)$,
and Cauchy--Schwarz gives $|x^*Ax|\le\|A\|$ for unit $x$.
The affine formula is an expansion; the two isometry formulas follow
by sending $x$ to $Ux$ or $Qx$, and using $U^*$ for the reverse inclusion.

For convexity, choose $x,y$ representing two points. Their complex
span has dimension at most two. In dimension one the compression
is scalar and the two points coincide. In dimension two choose an
orthonormal basis $Q$ of that span. Both points belong to $W(Q^*AQ)$;
the preceding lemma puts their segment there, and compression
inclusion puts the segment in $W(A)$. The empty case is convex too.

For (S11), $\operatorname{Re}(x^*Ax)=x^*Hx$. Diagonalizing $H$
expresses these values as $\sum_j\lambda_j|x_j|^2$ and bounds them
between the two extreme eigenvalues. Both endpoints occur.
Every intermediate value occurs by combining normalized eigenvectors
of the two extremes with coefficients $\sqrt t,\sqrt{1-t}$;
if the extremes coincide there is only one value.
For normal $A$, its unitary diagonalization similarly gives
$x^*Ax=\sum_j\lambda_j|x_j|^2$.
Conversely any finite convex combination of the eigenvalue list is
realized by $x_j=\sqrt{w_j}$. Repeated eigenvalues cause no problem;
the empty list gives the empty convex hull.
\end{proof}

\paragraph{Verified reuse and formalization boundary.}
\code{NLAlib.exists_unitary_diagonalization_of_normal} and
\code{NLAlib.normal_iff_exists_orthonormalBasis_eigenvectors} are
already indexed. Orthonormal completion and Hermitian spectral
decomposition are existing matrix interfaces. The new geometric
work is precisely (S10) and the explicit one-dimensional correction
inside $\ker L$; there is no contour integral or separation theorem.

\section{Positive Hermitian part: kernels and scalar products}
\label{struct:accretive}
\paragraph{Catalogue target.} \code{psd-hermitian-part}.
Write $A=H+iK$ with $H=H^*$ and $K=K^*$.

\begin{theorem}[Kernel, rank, and row/column inclusion]
If $H\succeq0$, then
\[
 \ker A=\ker A^*=\ker H\cap\ker K,
 \qquad \rank H\le\rank A.
 \tag{S12}
\]
If additionally $\rank H=\rank A$, $A$ has both row and column
inclusion for every leading principal cut. In particular this holds
when $H\succ0$, in which case every principal block of $A$ is invertible.
\end{theorem}
\begin{proof}
If $Ax=0$, its quadratic form has real part $x^*Hx=0$.
Diagonalizing the positive semidefinite $H$ shows $Hx=0$,
and $Ax=0$ then gives $Kx=0$. The converse is immediate, and the
same proof applies to $A^*=H-iK$. Rank--nullity gives the rank bound.
With equality of ranks, inclusion of kernels gives
$\ker A=\ker H$ and hence $\ker H\subseteq\ker K$.

Let $I$ be a leading index set, $J$ its complement, and
$B=A_{II}$. If $Bz=0$, extend $z$ by zero on $J$ to $\widetilde z$.
Then $\widetilde z^*H\widetilde z=\operatorname{Re}(z^*Bz)=0$,
so $H\widetilde z=0$, then $K\widetilde z=0$ and
$A\widetilde z=A^*\widetilde z=0$. In particular
$A_{JI}z=0$. Thus
\[
 \ker B\subseteq\ker A_{JI},
\]
which is equivalent to every row of $A_{JI}$ lying in the row span
of $B$: a linear functional vanishing on $\ker B$ factors through
$z\mapsto Bz$, then extends from $\ran B$ to the whole codomain.
Applying the same argument to $B^*$ gives
$\ker B^*\subseteq\ker A_{IJ}^*$, equivalent to every column of
$A_{IJ}$ lying in the column space of $B$.
This proves the required local row and column conditions, in fact
for all rows or columns beyond each cut.
If $H\succ0$, every principal $H_{II}$ is positive definite,
so $Bz=0$ forces $z=0$ without the rank-equality argument.
\end{proof}

Assume now $H\succ0$ and $n\ge1$, and let
\[
 T=H^{-1/2}KH^{-1/2}=T^*,\qquad t_1,\ldots,t_n\in\R
\]
be its eigenvalues. Congruence and the Hermitian spectral theorem give
the exact identities
\[
 |\det A|=\det H\prod_{j=1}^n\sqrt{1+t_j^2},\qquad
 |\det K|=\det H\prod_{j=1}^n|t_j|.
 \tag{S13}
\]
Only this simultaneous reduction is needed; $H$ and $K$ need not commute.

\begin{lemma}[Finite product inequality]
For $x_j\ge0$ and $n\ge1$,
\[
 \prod_{j=1}^n(1+x_j)\ge
   \left(1+\left(\prod_{j=1}^nx_j\right)^{1/n}\right)^n.
 \tag{S14}
\]
\end{lemma}
\begin{proof}
If some $x_j=0$, the right side is $1$ and the claim is immediate.
Otherwise expand the left side as $\sum_{r=0}^n e_r(x)$.
The $\binom nr$ terms in $e_r$ have product
$(\prod_jx_j)^{\binom{n-1}{r-1}}$, because each index occurs in that
many subsets. Arithmetic--geometric mean therefore gives
\[
 e_r(x)\ge\binom nr\left(\prod_jx_j\right)^{r/n}.
\]
Summing is the binomial formula. For $n>1$, equality holds exactly
when all $x_j$ agree: necessity follows already from the $r=1$
AM--GM inequality in the positive case, and from the product being
$1$ in the zero case. For $n=1$, equality is automatic.
\end{proof}

\begin{theorem}[Ostrowski--Taussky and Fan with exact constants]
If $H\succ0$, then
\[
 \det H\le|\det A|,
 \tag{S15}
\]
with equality exactly when $K=0$. For $n\ge1$,
\[
 (\det H)^{2/n}+|\det K|^{2/n}\le|\det A|^{2/n}.
 \tag{S16}
\]
For $n\ge2$,
\[
 \det H+|\det K|\le|\det A|.
 \tag{S17}
\]
\end{theorem}
\begin{proof}
Equation (S15) follows from (S13), since every factor is at least one;
equality forces every $t_j=0$, hence $T=K=0$.
Apply (S14) to $x_j=t_j^2$, take the $n$-th root, and multiply by
$(\det H)^{2/n}$ to obtain (S16).
For (S17), set $x_j=t_j^2$ and $v=\prod_j|t_j|$.
In the subset expansion of $\prod_j(1+x_j)$ retain the empty subset,
the full subset, $\{1\}$, and its complement. These are four distinct
subsets when $n\ge2$. All other summands are nonnegative, and
\[
 \prod_j(1+x_j)\ge1+v^2+x_1+\prod_{j=2}^n x_j
       \ge1+v^2+2v=(1+v)^2.
\]
The second inequality is just $a+b\ge2\sqrt{ab}$ for nonnegative
$a,b$. Taking square roots and multiplying by $\det H$ proves the
claim using only a finite expansion and a quadratic inequality.
\end{proof}

\paragraph{Boundary cases and equality audit.}
The restriction $n\ge2$ in (S17) is essential:
$H=K=[1]$ would assert $2\le\sqrt2$ in dimension one.
The exponents in (S16) require $n\ge1$; dimension zero is not an
instance of that formula. Equation (S15) does include dimension zero,
where both determinants are $1$ and $K$ is the zero matrix.
For $n>1$, equality in (S16) is equivalent to
$T^2=t^2I$ for some $t\ge0$. In (S17), equality for $n=2$ means
$|t_1|=|t_2|$, and for $n>2$ it means $K=0$.

\paragraph{Verified reuse.}
Hermitian spectral decomposition, positive square roots and congruence
are existing matrix interfaces. The currently indexed row-inclusion
theorems, including \code{NLAlib.exists_isLU_iff_rowInclusion}, use real
scalars; their row-annihilator proof extends over $\C$, but the existing
signature must not be cited as the complex theorem. The new work here
is (S12), the kernel-factorization argument at a principal cut, and the
finite scalar inequality (S14). No numerical-range convexity is needed.


\section{Orthogonal polynomials from finite moment systems}
\label{sec:orthogonal-foundations}
\paragraph{Targets and conventions.}
This section covers \code{orthogonal-polynomials-measure-def},
\code{hankel-moment-positivity}, and \code{formal-orthogonal-polynomials}.
Let $\mu$ be a finite positive Borel measure on $\R$, with all absolute
moments finite. For a result using polynomials of degree at most $d$, it is
enough to assume finiteness of the moments actually used (in particular,
the absolute moment of order $2d$ for their Gram matrix). Write
\[
 L(f)=\int f(t)\,d\mu(t),\quad m_r=L(t^r),\quad
 M_n=(m_{i+j})_{0\le i,j<n},\quad \Delta_n=\det M_n,\quad \Delta_0=1.
\]
The useful primitive is the finite bilinear form $L(fg)$ on polynomials of
bounded degree. No global inner product on $\R[X]$, no completion, and no
measure-dependent Hilbert-space instance is needed.

\begin{lemma}[The support interface]
For every real polynomial $p$ whose square is integrable,
\[
 L(p^2)=0\quad\Longleftrightarrow\quad
 p=0\text{ on }\operatorname{supp}\mu.
\]
Consequently $L(pq)$ is positive definite on the space $\mathcal P_d$ of
polynomials of degree at most $d$ if and only if the support contains at
least $d+1$ distinct points.
\end{lemma}
\begin{proof}
If $L(p^2)=0$, then $p^2=0$ almost everywhere. If $p(t_0)\ne0$ at a support
point, continuity gives a neighbourhood where $p^2$ is bounded below by a
positive constant, and that neighbourhood has positive measure, a
contradiction. Conversely a finite Borel measure on $\R$ gives full measure
to its support, so vanishing there implies zero integral. A nonzero
polynomial of degree at most $d$ has at most $d$ roots. If instead the
support has at most $d$ points, the product of $X-t$ over those points is a
nonzero polynomial in $\mathcal P_d$ of zero norm. The empty product is $1$;
thus this argument includes the zero measure.
\end{proof}

\paragraph{Existence by one linear solve.}
Assume that the support contains at least $k$ points, so $M_k$ is positive
definite and invertible. Put $v_k=(m_k,\ldots,m_{2k-1})^\top$ and solve
\[
 M_kc=-v_k,\qquad \pi_k(X)=X^k+\sum_{j<k}c_jX^j.                 \tag{OF1}
\]
For $k=0$ this means $\pi_0=1$. The equations say exactly
$L(\pi_kX^j)=0$ for $j<k$. They prove existence and uniqueness of the monic
degree-$k$ polynomial orthogonal to all lower degrees. Its norm
$h_k=L(\pi_k^2)$ is positive if the support has at least $k+1$ points.
The triangular coefficient matrix with diagonal $1$ shows that
$\pi_0,\ldots,\pi_d$ span $\mathcal P_d$. Orthogonality to all of them
therefore forces a member of $\mathcal P_d$ to vanish whenever that form
is positive definite.

If the support consists of exactly $N$ distinct points, the positive-norm
family is precisely $\pi_0,\ldots,\pi_{N-1}$, and (OF1) still uniquely
defines
\[
 \pi_N(X)=\prod_{t\in\operatorname{supp}\mu}(X-t),\qquad h_N=0.   \tag{OF2}
\]
Indeed the displayed polynomial is monic and has zero product integral
against every polynomial, so uniqueness in (OF1) applies. There is no
division by $h_N$ and no normalized $p_N$. Beyond this point the phrase
``the orthogonal polynomial'' requires an additional convention and is
not supplied by positive-definite orthogonalization. For $\mu=0$, $N=0$,
the positive-norm family is empty and the terminal polynomial is $1$.

\paragraph{Hankel positivity without a general Sylvester theorem.}
If $p=\sum_{j<n}c_jX^j$, finite linearity of the integral gives
\[
 L(p^2)=c^\top M_nc.                                           \tag{OF3}
\]
For a positive measure each $M_n$ is already PSD. Thus $M_n\succ0$ is
equivalent simply to $\Delta_n>0$: diagonalize $M_n$, whose eigenvalues
are nonnegative, and use their product. Moreover $M_n\succ0$ implies all
its leading principal submatrices are PD. It follows immediately that
\[
 M_n\succ0\quad\Longleftrightarrow\quad
 \Delta_1,\ldots,\Delta_n>0.
\]
This special-purpose proof avoids formalizing Sylvester's criterion for
arbitrary symmetric matrices. At $n=0$, both sides are vacuous truths.

\paragraph{Determinants and coefficients.}
Whenever $\Delta_k\ne0$, finite determinant expansion gives
\[
 \pi_k(X)=\frac1{\Delta_k}
 \det\begin{pmatrix}
 m_0&m_1&\cdots&m_k\\
 m_1&m_2&\cdots&m_{k+1}\\
 \vdots&\vdots&&\vdots\\
 m_{k-1}&m_k&\cdots&m_{2k-1}\\
 1&X&\cdots&X^k
 \end{pmatrix}.                                               \tag{OF4}
\]
The coefficient of $X^k$ is $1$; applying $L(X^j\cdot)$ for $j<k$
repeats one of the first $k$ rows, proving (OF1). The $k=0$ formula is
the $1\times1$ determinant $1$.
Define $E_0=0$, and for $k\ge1$ define $E_k$ as the determinant with row
indices $0,\ldots,k-1$ and column indices $0,\ldots,k-2,k$ of the Hankel
array $(m_{i+j})$. This explicit definition removes the ambiguity in a
``deleted penultimate column'' convention. If $c_k$ is the coefficient of
$X^{k-1}$ in $\pi_k$, (OF4) yields $c_k=-E_k/\Delta_k$.
The unit triangular change from monomials to monic orthogonal polynomials
gives
\[
 \Delta_n=\prod_{j<n}h_j,\qquad h_k=\frac{\Delta_{k+1}}{\Delta_k}.
                                                                    \tag{OF5}
\]
The recurrence derived below therefore has
\[
 \alpha_k=\frac{E_{k+1}}{\Delta_{k+1}}-\frac{E_k}{\Delta_k},\qquad
 \beta_k=\frac{h_k}{h_{k-1}}
        =\frac{\Delta_{k+1}\Delta_{k-1}}{\Delta_k^2}\quad(k\ge1). \tag{OF6}
\]
Only use these quotients when their denominators are nonzero. In the
finite-support case $\beta_N=0$ is valid, but $\alpha_N$ is not defined by
a norm quotient.

\paragraph{Spectral-measure bridge.}
For real symmetric $A$ and any $b$, let
$\mu_b=\sum_i(u_i^\top b)^2\delta_{\lambda_i}$, without normalizing its
mass. The spectral theorem gives
\[
 L(pq)=\langle p(A)b,q(A)b\rangle,\qquad
 M_n=K_n^\top K_n,\quad K_n=[b,Ab,\ldots,A^{n-1}b].             \tag{OF7}
\]
Consequently $M_n\succ0$ iff $n\le\operatorname{grade}(A,b)$.
If $b=0$, the grade is zero and only the empty Gram matrix is PD.
In the nonzero case the monic polynomial of order $k$ agrees with
$\det(XI-T_k)$: both are monic of degree $k$ and orthogonal to every lower
degree, including the terminal case $k=\operatorname{grade}(A,b)$.

\paragraph{Formal orthogonality is purely bilinear.}
For $L:\C[X]\to\C$ linear, use $B(p,q)=L(pq)$, not $L(p\overline q)$.
Assume every $\Delta_n\ne0$. The same systems (OF1) give unique monic
$\pi_k$ orthogonal to lower degrees; the unit triangular determinant
identity (OF5) gives $h_k=L(\pi_k^2)\ne0$. Hence division by $h_k$ is valid,
even though positivity and real zeros are unavailable. Conversely,
a monic orthogonal sequence with every $h_k\ne0$ gives
$\Delta_n=\prod_{k<n}h_k\ne0$. Alternatively, uniqueness of the monic
lower-degree orthogonality problem for every $k$ makes each square system
(OF1) uniquely solvable and therefore makes $M_k$ invertible.
Merely assuming pairwise orthogonality, without uniqueness or nonzero
norms, is insufficient: the zero functional is an immediate counterexample.
The multiplication symmetry $B(Xp,q)=B(p,Xq)$ gives the same three-term
recurrence below, now with $\beta_k\ne0$ rather than $\beta_k>0$.

\paragraph{Verified reuse and remaining interface.}
The declaration snapshot contains
\code{NLAlib.integral_eval_mul_eval_spectralMeasure},
\code{NLAlib.integral_pow_spectralMeasure},
\code{NLAlib.linearIndependent_pow_mulVec_iff_le_krylovGrade},
\code{NLAlib.LanczosDecomp.charpoly_T_monic}, and
\code{NLAlib.LanczosDecomp.integral_eval_charpoly_mul_eq_zero}.
They directly supply (OF7), the grade test, and the final uniqueness bridge.
For a general measure the new interface is finite moment integrability,
the support lemma, and the coefficient-system definition (OF1).

\section{The Jacobi matrix is compressed multiplication}
\label{sec:monic-recurrence}
\paragraph{Target and normalization.}
The target is \code{monic-orthogonal-recurrence}. Whenever $h_j>0$ set
\[
 p_j=\pi_j/\sqrt{h_j},\quad h_0=\mu(\R),\quad
 \alpha_j=\frac{L(X\pi_j^2)}{h_j},\quad
 \beta_j=\frac{h_j}{h_{j-1}},\quad a_j=\sqrt{\beta_j}\quad(j\ge1).
\]
Here $\beta_j$ is the \emph{monic} recurrence coefficient, whereas
\code{NLAlib.LanczosDecomp.beta} denotes the off-diagonal $a_j$.
Keeping these separate prevents a missing square in every norm or
continued-fraction formula. Put $\pi_{-1}=0$, $\pi_0=1$; no actual negative
index is needed in a formal implementation.

\begin{theorem}[Three-term recurrence and the norm product]
Below the positive-norm termination index,
\[
 \pi_{j+1}=(X-\alpha_j)\pi_j-\beta_j\pi_{j-1},\qquad
 h_j=h_0\prod_{i=1}^{j}\beta_i=h_0\prod_{i=1}^{j}a_i^2.         \tag{OF8}
\]
The first recurrence is $\pi_1=X-\alpha_0$.
It also constructs the terminal $\pi_N$ from $\pi_{N-1},\pi_{N-2}$
when the support has $N$ points.
\end{theorem}
\begin{proof}
Expand $X\pi_j-\pi_{j+1}$ in the monic basis
$\pi_0,\ldots,\pi_j$. For $i\le j-2$ its coefficient times $h_i$ is
$L(X\pi_j\pi_i)=L(\pi_jX\pi_i)=0$ because $\deg(X\pi_i)<j$.
The coefficient at $\pi_j$ is $\alpha_j$.
For $j\ge1$, $X\pi_{j-1}-\pi_j$ has degree below $j$, so the coefficient
at $\pi_{j-1}$ times $h_{j-1}$ is $L(\pi_jX\pi_{j-1})=h_j$.
This proves the recurrence; telescoping the norm ratios proves the
product. The proof only divides by norms before the terminal index.
\end{proof}

In the orthonormal basis, compression of multiplication by $X$ to
$\mathcal P_{q-1}$ is the symmetric matrix
\[
 J_q=\begin{pmatrix}
 \alpha_0&a_1&&\\a_1&\alpha_1&\ddots&\\
 &\ddots&\ddots&a_{q-1}\\&&a_{q-1}&\alpha_{q-1}
 \end{pmatrix}.
\]
Expansion of $\det(XI-J_{j+1})$ in its last row and then the remaining
last column gives
$D_{j+1}=(X-\alpha_j)D_j-a_j^2D_{j-1}$, with $D_0=1,D_1=X-\alpha_0$.
Induction against (OF8) proves $D_j=\pi_j$. This finite determinant
identity is also the denominator recurrence for the convergents of the
Jacobi continued fraction.

\paragraph{Christoffel--Darboux without division by $x-y$.}
For $m\ge1$ with $h_0,\ldots,h_{m-1}>0$, multiplication of the two
recurrences by $\pi_j(y)$ and $\pi_j(x)$ and subtraction gives a
telescoping identity:
\[
 (x-y)\sum_{j<m}\frac{\pi_j(x)\pi_j(y)}{h_j}
 =\frac{\pi_m(x)\pi_{m-1}(y)-\pi_{m-1}(x)\pi_m(y)}{h_{m-1}}.
                                                                    \tag{OF9}
\]
This is a polynomial identity for \emph{all} $x,y$, including $x=y$;
there is no removable-singularity interface to formalize. Multiplying by
$h_n$ and setting $m=n+1$ gives the monic formula
\[
 (x-y)\sum_{k=0}^{n}\left(\prod_{i=k+1}^{n}\beta_i\right)
 \pi_k(x)\pi_k(y)
 =\pi_{n+1}(x)\pi_n(y)-\pi_n(x)\pi_{n+1}(y).                  \tag{OF10}
\]
Differentiate (OF9) formally in $x$ and then substitute $y=x$:
\[
 \sum_{j<m}p_j(x)^2
 =\frac{\pi_m'(x)\pi_{m-1}(x)-\pi_{m-1}'(x)\pi_m(x)}{h_{m-1}}
 >0.                                                          \tag{OF11}
\]
Strict positivity follows from the term $p_0(x)^2=1/h_0$.
If also $h_m>0$, the right side is
$a_m[p_m'(x)p_{m-1}(x)-p_{m-1}'(x)p_m(x)]$.
At finite-support termination use (OF11), not this normalized version,
because $p_m$ then does not exist.

\paragraph{Verified reuse.}
\code{NLAlib.LanczosDecomp.X_mul_poly},
\code{NLAlib.LanczosDecomp.charpoly_restrict_T_mul_poly_zero},
\code{NLAlib.LanczosDecomp.integral_poly_mul_poly}, and
\code{NLAlib.LanczosDecomp.sub_mul_sum_eval_poly_mul}
already provide the normalized recurrence, normalization bridge,
orthogonality, and nonconfluent identity in their stated index ranges.
The additional general-measure route above uses finite basis expansion
and formal polynomial differentiation; it does not require differentiating
an integral.

\paragraph{Terminal convention.}
At $N$-point termination with $N\ge1$, define
$\beta_N=h_N/h_{N-1}=0$ and $a_N=0$.
Then the norm product (OF8) remains valid at $j=N$, with both sides zero;
it does not define $p_N$. The characteristic-polynomial bridge (OF7)
can be justified either by the checked orthogonality declaration or by
the determinant recurrence just proved, so it is not an extra assumption.

\section{Zeros and strict interlacing from an unreduced tridiagonal matrix}
\label{sec:orthogonal-zeros}
\paragraph{Targets.}
This covers \code{orthogonal-polynomial-zeros} and the strict-interlacing
part of \code{jacobi-matrix-identities}. Suppose the support has at least
$q$ points and $q\ge1$, so $h_0,\ldots,h_{q-1}>0$ and
$a_1,\ldots,a_{q-1}>0$.

\begin{lemma}[One coordinate determines every Jacobi eigenvector]
For an eigenvector $v$ of $J_q$, its first coordinate determines all
coordinates. Every nonzero eigenvector has $v_0\ne0$, every eigenvalue
is simple, and every squared first coordinate is positive.
The same conclusion holds for the last coordinate.
\end{lemma}
\begin{proof}
The first row gives $a_1v_1=(\theta-\alpha_0)v_0$; each successive row
gives $a_{j+1}v_{j+1}=(\theta-\alpha_j)v_j-a_jv_{j-1}$.
All coefficients used to divide are positive. An eigenvector with
$v_0=0$ vanishes by induction. Evaluation at the first coordinate is an
injective linear map from the eigenspace into $\R$, so its dimension is
at most one. A real symmetric matrix has an orthonormal eigenbasis;
therefore algebraic and geometric multiplicity agree and each eigenvalue
is simple. Reverse the coordinate order for the final claim. For $q=1$
the conclusions are immediate without a recurrence step.
\end{proof}

Thus $\pi_q=\det(XI-J_q)$ has $q$ distinct real roots. If the support has
at least $q+1$ points, $a_q>0$ and $\pi_q,\pi_{q+1}$ have no common
root: a common root in (OF8) would successively force
$\pi_{q-1},\ldots,\pi_0$ to vanish there, contradicting $\pi_0=1$.
Cauchy's weak interlacing for the leading principal submatrix $J_q$ of
$J_{q+1}$ is already available. Equality in either inequality would be a
common root, so with eigenvalues in decreasing order,
\[
 \theta_{i+1}(J_{q+1})<\theta_i(J_q)<\theta_i(J_{q+1})
 \qquad(1\le i\le q).                                        \tag{OF12}
\]
This proof needs no confluent Christoffel--Darboux theorem and no
perturbation argument.

\paragraph{Strict location and spectral gaps.}
Suppose $\operatorname{supp}\mu\subseteq[a,b]$ and it has more than $q$
points. For a unit eigenvector $v$ put $r=\sum_{j<q}v_jp_j$.
Compression of multiplication gives
\[
 \theta=L(Xr^2),\quad L(r^2)=1,\quad
 \theta-a=L((X-a)r^2),\quad b-\theta=L((b-X)r^2).
\]
Both last quantities are positive. At most one support point is $a$,
whereas the nonzero polynomial $r$ has at most $q-1$ roots; hence there
is a support point strictly above $a$ where $r\ne0$. Continuity and the
support interface make the first integral positive. The other endpoint
is identical. Thus $a<\theta<b$. With an unbounded support, use every
finite lower or upper support bound that exists; this gives the usual
open-convex-hull statement. At finite-support termination the roots are
the support points, so strict endpoint inequalities must not be asserted.

For the gap claim, let $(c,d)$ contain no support point and suppose two
distinct zeros $r<s$ of $\pi_q$ lie there. If the support has more than
$q$ points, put $g=\pi_q/((X-r)(X-s))$, a polynomial of degree $q-2$.
Orthogonality gives
\[
 0=L(\pi_qg)=L\bigl((X-r)(X-s)g^2\bigr).                     \tag{OF13}
\]
On the support $(t-r)(t-s)>0$. The polynomial $g$ cannot vanish on the
whole support, so the last integral is strictly positive, a
contradiction. This proof only integrates a polynomial. The claim is
automatic for $q=1$ and, at $q=N$, from the node polynomial (OF2).

\paragraph{Checked reuse.}
% Exact names: NLAlib.LanczosDecomp.eigenvalues₀_restrict_T_le
% and NLAlib.LanczosDecomp.eigenvalues₀_le_eigenvalues₀_restrict_T.
\code{NLAlib.LanczosDecomp.eigenvalues}\textsubscript{0}\code{_restrict_T_le} and
\code{NLAlib.LanczosDecomp.eigenvalues}\textsubscript{0}\code{_le_eigenvalues}\textsubscript{0}\code{_restrict_T}
provide weak interlacing. The snapshot also contains
\code{NLAlib.LanczosDecomp.beta_ne_zero} for arbitrary valid
decompositions, and
\code{NLAlib.LanczosDecomp.beta_arnoldiDecomp_pos} for the canonical
construction. Nonzero off-diagonals suffice for simplicity and strict
interlacing; positive signs implement our normalized polynomial convention.

\section{Christoffel numbers from a finite orthogonal matrix}
\label{sec:christoffel-numbers}
The target is \code{gauss-christoffel-numbers}. Let $q\ge1$ and let
$\theta_i$ be the distinct eigenvalues of $J_q$ with unit eigenvectors
$u_i$. Evaluate the recurrence at a zero of $\pi_q$ to obtain
\[
 J_q\mathbf p(\theta_i)=\theta_i\mathbf p(\theta_i),\qquad
 \mathbf p=(p_0,\ldots,p_{q-1})^\top.                          \tag{OF14}
\]
In the last row the discarded term is proportional to $\pi_q$, so this
remains valid at termination where $p_q$ is undefined. Since
$p_0=1/\sqrt{h_0}$, the vector is nonzero. The one-dimensional eigenspace
therefore gives
\[
 u_i=\pm\frac{\mathbf p(\theta_i)}{\|\mathbf p(\theta_i)\|},\qquad
 w_i=h_0(u_i)_0^2
     =\left(\sum_{j<q}p_j(\theta_i)^2\right)^{-1}>0.            \tag{OF15}
\]
The $w_i$ are weights for $\mu$, whereas $(u_i)_0^2$ are weights for
$\mu/h_0$. For $\mu_b$ the mass is $h_0=\|b\|^2$.
Put the eigenvectors in the columns of an orthogonal matrix $U$.
The identity $UU^\top=I$ gives directly
\[
 \sum_iw_ip_j(\theta_i)p_k(\theta_i)=\delta_{jk}\quad(j,k<q),
 \qquad\sum_iw_i=h_0.                                       \tag{OF16}
\]
Thus discrete orthogonality does not need a pre-existing Gauss rule.
Conversely, for $\deg f\le2q-1$ divide $f=s\pi_q+r$, where
$\deg s<q$, $\deg r<q$. The first part has both zero integral and zero
nodal values. Expanding $r$ in the $p_j$ and using (OF16) with $k=0$ yields
\[
 L(f)=\sum_iw_if(\theta_i)=h_0e_0^\top f(J_q)e_0.              \tag{OF17}
\]
The empty rule $q=0$ has mass zero; it cannot represent a nonzero measure.

\paragraph{Checked reuse.}
\code{NLAlib.LanczosDecomp.cfc_T_apply_zero_zero},
\code{NLAlib.LanczosDecomp.dotProduct_cfc_T_mulVec_eq_sum}, and
\code{NLAlib.LanczosDecomp.dotProduct_aeval_mulVec_eq_sum}
already give the spectral sum and polynomial quadrature for a Lanczos
decomposition. The new short interface is (OF14) and its normalization.

\section{Recurrence-coefficient bounds by scalar Cauchy--Schwarz}
\label{sec:coefficient-bounds}
The target is \code{recurrence-coefficient-bounds}. Suppose
$\operatorname{supp}\mu\subseteq[a,b]$ and it has more than $k+1$ points.
The probability measure $p_k^2\mu$ gives
\[
 \alpha_k=L(Xp_k^2),\qquad a<\alpha_k<b.                      \tag{OF18}
\]
For strictness $(X-a)p_k^2$ and $(b-X)p_k^2$ cannot vanish identically
on the support: outside either endpoint there are more than $k$ support
points, but $p_k$ has degree $k$.
For $k\ge1$ with $h_k>0$, Cauchy--Schwarz gives
\[
 0<a_k=L(Xp_{k-1}p_k)
 \le\sqrt{L(X^2p_{k-1}^2)}\le\max\{|a|,|b|\},
 \qquad 0<\beta_k=a_k^2\le\max\{a^2,b^2\}.                   \tag{OF19}
\]
There is no $k=0$ norm-ratio coefficient here. Since
$L(p_{k-1}p_k)=0$, replacing $X$ by $X-(a+b)/2$ also proves the optional
stronger bound $a_k\le(b-a)/2$.
For a symmetric matrix, the same finite-vector calculations are
$\alpha_j=v_j^\top Av_j$ and $a_j=v_{j-1}^\top Av_j$, with orthonormal
$v_j$. Rayleigh's inequality and Cauchy--Schwarz give
\[
 \lambda_{\min}(A)\le\alpha_j\le\lambda_{\max}(A),\qquad
 0<a_j\le\|A\|_2\quad(1\le j<\operatorname{grade}(A,b)).
\]
The next off-diagonal at termination is zero. The catalogue uses $\beta$
for both $a_j$ and $a_j^2$ in different clauses; (OF19) distinguishes them.

\section{Discrete measures as diagonal matrix problems}
\label{sec:discrete-jacobi}
The target is \code{discrete-measure-jacobi-matrix}. Let
$\mu=\sum_{r=1}^Nw_r\delta_{t_r}$, with distinct real $t_r$ and $w_r>0$.
Put $D=\diag(t_1,\ldots,t_N)$ and $b_r=\sqrt{w_r}$. Then
\[
 L(f)=b^\top f(D)b,\qquad
 K_N=\diag(\sqrt w)\,(t_r^{\,j})_{r,\,0\le j<N}.              \tag{OF20}
\]
The Vandermonde determinant and positive weights make $K_N$ invertible,
so the grade is $N$. This proves the exact dimension without analyzing
a recurrence algorithm. For $N\ge1$, define
$Q_{rj}=\sqrt{w_r}\,p_j(t_r)$, $0\le j<N$.
Orthogonality gives $Q^\top Q=I$; as $Q$ is square, $QQ^\top=I$ too.
The recurrence and the terminal node polynomial give
\[
 Q^\top DQ=J_N,\qquad Q^\top b=\sqrt{h_0}\,e_0,\qquad
 \sum_{j<N}p_j(t_r)p_j(t_s)=\frac{\delta_{rs}}{w_r}.            \tag{OF21}
\]
For the last formula $QQ^\top=I$ first gives denominator
$\sqrt{w_rw_s}$; this equals $w_r$ on the diagonal and is irrelevant
when $r\ne s$ because the numerator is zero. The coefficients of the
discrete monic polynomials and those of a positively oriented graded
compression of $(D,b)$ therefore coincide, by uniqueness.
This is an identity of mathematical constructions, with no procedural
claim. For $N=0$ the matrices are empty and the grade is zero; no
first-coordinate vector $e_0$ is invoked.

Conversely the spectral measure of $(J_q,\sqrt{h_0}e_0)$ is exactly
$\sum_iw_i\delta_{\theta_i}$ by (OF15). For matrix spectral measures take
$h_0=\|b\|^2$. The definition \code{NLAlib.spectralMeasure} and the
checked integral and grade identities in (OF7) are existing entry points;
the new general discrete bridge is the explicit evaluation matrix $Q$.

\section{Symmetry and squaring the variable}
\label{sec:symmetric-measure}
The target is \code{symmetric-measure-even-odd-reduction}.
If reflection $t\mapsto-t$ preserves $\mu$, the polynomial
$(-1)^k\pi_k(-X)$ is monic of degree $k$ and orthogonal to lower degrees,
by change of variables. Uniqueness gives
\[
 \pi_k(-X)=(-1)^k\pi_k(X),\qquad \alpha_k=0\quad(h_k>0).       \tag{OF22}
\]
The second assertion holds because $X\pi_k^2$ is odd. The parity assertion
also includes the terminal polynomial of a finite-support measure.
Let $s(t)=t^2$, $\nu_+=s_*\mu$, and $\nu_-=s_*(t^2\mu)$, where
$\int f\,d(t^2\mu)=\int t^2f(t)\,d\mu$. Parity gives monic $r_k^+,r_k^-$
such that
\[
 \pi_{2k}(t)=r_k^+(t^2),\qquad
 \pi_{2k+1}(t)=t\,r_k^-(t^2).                                \tag{OF23}
\]
For $j<k$, the two orthogonality calculations are
\[
 \int r_k^+(u)u^j\,d\nu_+(u)=L(\pi_{2k}X^{2j})=0,\qquad
 \int r_k^-(u)u^j\,d\nu_-(u)=L(\pi_{2k+1}X^{2j+1})=0.
\]
The squared norms are $h_{2k}$ and $h_{2k+1}$. Thus these are precisely
the monic orthogonal polynomials for the push-forward measures, wherever
the positive-norm families exist, and with the same terminal interpretation
when a displayed norm vanishes. In particular $\nu_-$ can be zero (if
$\mu$ is supported at $0$); do not normalize it in that case.

\paragraph{The augmented matrix identity.}
For a real rectangular matrix $B$ put
\[
 M=\begin{pmatrix}0&B\\B^\top&0\end{pmatrix},\quad c=(u,0),\quad
 S=\diag(I,-I).
\]
Then $SMS=-M$ and $Sc=c$. Spectral projectors for opposite eigenvalues
are carried into one another by $S$, with equal squared projections of
$c$; hence the spectral measure of $(M,c)$ is symmetric. Moreover
\[
 M^{2j}c=((BB^\top)^ju,0),\qquad
 M^{2j+1}c=(0,B^\top(BB^\top)^ju).
\]
Taking moments gives
\[
 \nu_+=\mu_u(BB^\top),\qquad
 \nu_-=\mu_{B^\top u}(B^\top B).                              \tag{OF24}
\]
One can avoid spectral-projector bookkeeping altogether: the displayed
matrix-power identities prove equality of polynomial integrals, and
Lagrange interpolation on the union of two finite supports then proves
equality of their discrete measures, atom by atom. Odd moments similarly
prove the reflection identity. This also handles $u=0$ or $B^\top u=0$.

\section{Complex Jacobi resolvents by finite block elimination}
\label{sec:complex-jacobi-resolvent}
This is the complex-variable part of \code{jacobi-matrix-identities}.
Let $q\ge1$ and regard $J_q$ as complex Hermitian. For $z\notin\R$, every
real symmetric principal block has $zI-J$ invertible: the imaginary part
of $v^*(zI-J)v$ is $(\operatorname{Im}z)\|v\|^2$, so its kernel is zero.
Consequently all inverses and denominators in
\[
 e_0^*(zI-J_q)^{-1}e_0
 =\cfrac1{z-\alpha_0-
       \cfrac{a_1^2}{z-\alpha_1-
       \cfrac{a_2^2}{\ddots-
       \cfrac{a_{q-1}^2}{z-\alpha_{q-1}}}}}                   \tag{OF25}
\]
exist in the ordinary field sense. Formally use recursion: if $R_j$ is
the first diagonal resolvent entry of the trailing block on indices
$j,\ldots,q-1$, solving its two block rows gives
\[
 R_{q-1}=(z-\alpha_{q-1})^{-1},\qquad
 R_j=(z-\alpha_j-a_{j+1}^2R_{j+1})^{-1}.
\]
The scalar Schur-complement denominator is nonzero because the whole
trailing block and its next trailing block are both invertible.

For a more robust division interface put
$D_j(z)=\det(zI-J_q[\{j,\ldots,q-1\}])$, with $D_q=1$ and
$D_{q-1}=z-\alpha_{q-1}$. Laplace expansion and Cramer's rule give
\[
 D_j=(z-\alpha_j)D_{j+1}-a_{j+1}^2D_{j+2}\quad(j\le q-2),
 \qquad e_0^*(zI-J_q)^{-1}e_0=\frac{D_1(z)}{D_0(z)}.           \tag{OF26}
\]
The quotient is valid for every complex $z$ with $D_0(z)\ne0$, including
real points where an intermediate denominator in (OF25) vanishes.
There (OF25) is an identity of rational functions, not a literal nested
formula using totalized division. For instance at $z=0$ the matrix
$J_2=\left(\begin{smallmatrix}0&1\\1&0\end{smallmatrix}\right)$ has an
invertible $zI-J_2$, but the inner denominator in $1/(z-1/z)$ is zero.
The determinant quotient correctly gives zero.

The spectral decomposition also gives
\[
 h_0e_0^*(zI-J_q)^{-1}e_0=\sum_i\frac{w_i}{z-\theta_i}.        \tag{OF27}
\]
Together with (OF17), this is the $q$-point Gauss or Jacobi convergent of
$\int(z-t)^{-1}\,d\mu(t)$: its coefficients at infinity match moments
through order $2q-1$. For compact support, use the uniformly convergent
geometric series when $|z|$ exceeds the support radius. For a general
finite-moment measure one may instead assert equality of those finite
moment coefficients, without claiming convergence of an infinite power
series. The finite matrix identity itself requires no analytic
continued-fraction convergence theorem. At $q=0$, define the zero scalar
approximant separately instead of using a nonexistent first coordinate.

\paragraph{Checked reuse and new interfaces.}
\code{NLAlib.LanczosDecomp.inv_smul_one_sub_T_apply_zero_zero}
exists for real $x$ above the original spectrum. The additional proof is
nonreal invertibility plus finite Schur-complement recursion. The rational
Cramer formula precisely handles real intermediate poles. Strict
interlacing was handled independently in (OF12); the checked
Golub--Welsch and nonconfluent CD statements need no new proofs.


% Advanced orthogonal-polynomial and quadrature handoff; snapshot e5725bb.
% Mathematical derivations only.  No new Lean proof is claimed.

\section{Exactness and uniqueness of quadrature by polynomial division}
\label{sec:quadrature-exactness}
\textbf{Atlas item:} \code{quadrature-exactness-characterisation}.
Let $L(p)=\int p\,d\mu$, let $n\ge1$, and let
$\tau_1,\ldots,\tau_n$ be distinct real nodes.  Write
$\omega(t)=\prod_{i=1}^n(t-\tau_i)$ and let $\ell_i$ be their Lagrange
basis.  Only the moments used below are required to be finite.
For $0\le k\le n$, the rule $Q(f)=\sum_i\lambda_i f(\tau_i)$ is exact
on $\mathcal P_{n-1+k}$ if and only if
\begin{equation}
 \lambda_i=L(\ell_i)\quad(1\le i\le n),\qquad
 L(\omega s)=0\quad(s\in\mathcal P_{k-1}),
 \tag{AQ1}
\end{equation}
where $\mathcal P_{-1}=\{0\}$.
Indeed, necessity follows by testing $\ell_i$ and $\omega s$.
For sufficiency divide $p=\omega s+r$, with $\deg r<n$ and
$\deg s\le k-1$; then $Q(p)=Q(r)=L(r)=L(p)$.
This is one Euclidean-division lemma, with no quadrature-specific analysis.

Suppose $L(p^2)>0$ for every nonzero $p\in\mathcal P_n$; for a positive
measure this follows from more than $n$ support points.  No $n$-node
rule is exact in degree $2n$, since $Q(\omega^2)=0<L(\omega^2)$.
Exactness through degree $2n-1$ forces $\omega=\pi_n$, the unique monic
orthogonal polynomial, by (AQ1).  Conversely its roots and the
interpolatory weights give that exactness.  Moreover
\begin{equation}
 \lambda_i=L(\ell_i^2)>0,
 \tag{AQ2}
\end{equation}
because $Q(\ell_i^2)=\lambda_i$ and $\deg\ell_i^2\le2n-2$.
Thus the Gauss rule is the unique $n$-atom positive measure with the
first $2n$ moments of $\mu$, with uniqueness understood up to permutation.
The same reasoning even excludes any rule with fewer than $n$ nodes
matching these moments: square its nodal polynomial.

At exactly $n$ support points, the $n$-point Gauss rule is the measure
itself and is exact for every integrable function.  Its uniqueness from
the first $2n$ moments still follows from positive definiteness on
$\mathcal P_{n-1}$ and (AQ1); the assertion that degree $2n$ is impossible
does not hold in this terminal case.  The zero-node rule is exact on
$\mathcal P_{-1}$, and is exact on constants only for the zero measure.

For a symmetric matrix and a nonzero vector $b$, take
$L(p)=b^\top p(A)b$.  Existing Lanczos moment matching and characteristic
polynomial orthogonality identify $\pi_q=\det(tI-T_q)$, so (AQ1) gives
the intrinsic uniqueness of its spectral quadrature rule, without
re-proving Lanczos moment identities.
\textbf{Verified reuse:}
\code{NLAlib.LanczosDecomp.dotProduct_aeval_mulVec_eq_sum},
\code{NLAlib.LanczosDecomp.integral_eval_charpoly_mul_eq_zero}, and
\code{NLAlib.integral_eval_spectralMeasure}.
The general-measure theorem still needs the short division-and-integration
wrapper; the existing Lebesgue-interval specialization
\code{NLAlib.gauss_lebesgue_quadrature_of_orthogonal_nodal} is not that wrapper.

\section{Derivative remainders from an explicit finite Hermite interpolant}
\label{sec:gauss-derivative-remainder}
\textbf{Atlas item:} \code{gauss-quadrature-derivative-remainder}.
Assume $a<c$, $n\ge1$, $\mu$ is finite, positive and supported in $[a,c]$,
and its $n$-point Gauss rule exists.  Let $f\in C^{2n}([a,c])$.
There is no need to develop a general Hermite-interpolation package first.
With the Lagrange basis of the Gauss nodes define explicitly
\begin{align}
 H(t)=\sum_i\bigl[&f(\tau_i)
       (1-2\ell_i'(\tau_i)(t-\tau_i))\ell_i(t)^2\\
       &+f'(\tau_i)(t-\tau_i)\ell_i(t)^2\bigr].
 \tag{AQ3}
\end{align}
Direct evaluation proves $\deg H\le2n-1$,
$H(\tau_i)=f(\tau_i)$ and $H'(\tau_i)=f'(\tau_i)$.
Exactness gives $Q(f)=L(H)$.

For $x$ not a node, put
$C_x=(f(x)-H(x))/\pi_n(x)^2$.  The function
$g(t)=f(t)-H(t)-C_x\pi_n(t)^2$ vanishes at $x$ and has a double zero
at each node.  Repeated Rolle's theorem gives a zero of $g^{(2n)}$
between the extreme zeros.  Therefore
\begin{equation}
 f(x)-H(x)=\frac{f^{(2n)}(\xi_x)}{(2n)!}\pi_n(x)^2,
 \qquad \xi_x\in(a,c).
 \tag{AQ4}
\end{equation}
At nodes the remainder is zero.  The only new real-analysis helper is:
a $C^r$ function with zeros whose multiplicities sum to at least $r+1$
has a zero of its $r$-th derivative between the extreme distinct zeros.
Its proof is induction: differentiation retains multiplicity $m-1$
at each old zero and supplies one zero in each gap by Rolle.
No choice of a measurable map $x\mapsto\xi_x$ is used.

Writing $h_n=L(\pi_n^2)$, pointwise upper and lower bounds in (AQ4)
and positivity of $L$ prove
\begin{equation}
 \frac{\min_{[a,c]}f^{(2n)}}{(2n)!}h_n
 \le L(f)-Q(f)\le
 \frac{\max_{[a,c]}f^{(2n)}}{(2n)!}h_n.
 \tag{AQ5}
\end{equation}
If $h_n>0$, put $d=(2n)![L(f)-Q(f)]/h_n$.
Continuity and the intermediate-value theorem imply that
$f^{(2n)}(\xi)=d$ for some $\xi\in(a,c)$: otherwise its difference
from $d$ has one strict sign throughout the connected open interval.
(AQ4) then has that strict sign off the nodes, whose complement has
positive $\pi_n^2\mu$-mass, contradicting the definition of $d$.
This argument also handles the possibility that a minimum or maximum
of the derivative is attained only at an endpoint.
If $h_n=0$, the support is contained in the nodes, the rule is exact,
and any $\xi\in(a,c)$ works.  Thus
\begin{equation}
 L(f)-Q(f)=\frac{f^{(2n)}(\xi)}{(2n)!}\,h_n.
 \tag{AQ6}
\end{equation}
For a singleton interval retain exactness, but do not assert existence
of a point in the empty interval $(a,c)$.

For $L=L_b$, the already proved identity
\code{NLAlib.integral_eval_mul_eval_spectralMeasure} gives
$h_q=\|\pi_q(A)b\|^2$.  This proves the matrix derivative remainder.
On a positive spectral interval, $1/t$ and $1/t^2$ have positive
even derivatives, whereas $\log t$ has negative even derivatives.
Consequently Gauss gives lower bounds for the first two and an upper
bound for the logarithm.  If $b=0$, both sides are zero separately.
\textbf{Remaining interface:} the finite multiplicity version of Rolle;
ordinary Lagrange remainders alone do not justify doubled nodes.

\section{Radau and Lobatto: weighted moments first, Jacobi borders second}
\label{sec:radau-lobatto}
\textbf{Atlas item:} \code{gauss-radau-lobatto}.
Use monic recurrences
$\pi_{j+1}=(t-\alpha_j)\pi_j-\beta_j\pi_{j-1}$,
$\beta_j=h_j/h_{j-1}$, and symmetric Jacobi off-diagonals
$a_j=\sqrt{\beta_j}$.  Thus the catalogue's symmetric coupling
written $\beta_q$ is $a_q$ in this document.

\subsection*{The rules and their exactness}
For left Radau, let $\phi_q^a$ be the monic degree-$q$ polynomial
orthogonal for $d\nu_a=(t-a)d\mu$.  Use the nodes $a$ and the zeros of
$\phi_q^a$, and interpolatory weights for $\mu$.
Its nodal polynomial $(t-a)\phi_q^a$ is orthogonal for $\mu$ to
$\mathcal P_{q-1}$, so (AQ1), with $n=q+1$ and $k=q$, proves
exactness through degree $2q$.  Right Radau is identical using
$d\nu_c=(c-t)d\mu$ and fixed node $c$.
The moment matrix for each modified measure on $\mathcal P_{q-1}$
must be positive definite to specify its degree-$q$ monic polynomial;
zero norm at degree $q$ is allowed as a terminal case.
For example, at least $q+1$ support points of $\mu$ suffice for
the required left or right degree-$q$ construction, even if one
endpoint is an atom.

For Lobatto use
$d\nu_{ac}=(t-a)(c-t)d\mu$, its degree-$q$ monic polynomial
$\phi_q^{ac}$, and the nodes $a,c$ and its $q$ roots.
Then $(t-a)(t-c)\phi_q^{ac}$ is orthogonal for $\mu$ to
$\mathcal P_{q-1}$, so the $q+2$ node rule is exact through degree
$2q+1$.  Its interior weights are positive by testing
$(t-a)(c-t)r_i(t)^2$, where $r_i$ is the Lagrange basis polynomial
on the $q$ interior nodes.  Endpoint weights are obtained by testing
$(c-t)(\phi_q^{ac})^2$ and $(t-a)(\phi_q^{ac})^2$; they are
nonnegative, and strictly positive precisely when the corresponding
integrals are positive.

\textbf{Support correction.}  The precise construction requires positive
definiteness of the modified moment matrix on $\mathcal P_{q-1}$
and the indicated nonzero endpoint-weight integrals if all $q+2$
weights are to be positive.  A clean sufficient condition is at least
$q+2$ support points of $\mu$.  The catalogue assumption
$\operatorname{grade}(A,b)\ge q+1$ alone does not guarantee this:
if the measure has exactly $q+1$ atoms including both endpoints,
then $\nu_{ac}$ has only $q-1$ atoms.  A degree-$q$ interior
orthogonal polynomial is then not uniquely determined.
At such finite termination use the actual atomic measure, which is
already an exact quadrature rule for every function, instead of
asserting a nondegenerate $q+2$ node Lobatto rule.
For $q=0$, Lobatto is simply the two-endpoint linear rule,
with weights $L((c-t)/(c-a))$ and $L((t-a)/(c-a))$;
one may vanish.  Radau at $q=0$ is the one prescribed endpoint
with weight $L(1)$.

\subsection*{Explicit borders and their moment proof}
Let $J_q$ be the original symmetric Jacobi matrix.
For $q\ge1$, left Radau is represented by
\begin{equation}
 J^R_{q+1}=\begin{pmatrix}J_q&a_qe_q\\a_qe_q^\top&\alpha^R\end{pmatrix},
 \qquad
 \alpha^R=a-\beta_q\frac{\pi_{q-1}(a)}{\pi_q(a)}.
 \tag{AQ7}
\end{equation}
The characteristic-polynomial expansion is
$(t-\alpha^R)\pi_q(t)-\beta_q\pi_{q-1}(t)$, which vanishes at $a$.
It is orthogonal to $\mathcal P_{q-1}$; dividing by $t-a$ identifies
it with $\phi_q^a$ by uniqueness.  The denominators are nonzero when
the original measure has more than $q$ support points, since the
zeros of $\pi_q$ lie strictly between the extreme support points.
Right Radau substitutes $c$ for $a$.

Under the sufficient nondegeneracy condition above, let $J_{q+1}$
be the original Jacobi block, set
$r(x)=\pi_q(x)/\pi_{q+1}(x)$, and define
\begin{equation}
 \gamma=\frac{c-a}{r(c)-r(a)}>0,
 \qquad d=a-\gamma r(a)=c-\gamma r(c),
 \qquad
 J^L_{q+2}=\begin{pmatrix}J_{q+1}&\sqrt\gamma e_{q+1}\\
                  \sqrt\gamma e_{q+1}^{\top}&d\end{pmatrix}.
 \tag{AQ8}
\end{equation}
The positivity follows from
$r(x)=e_{q+1}^{\top}(xI-J_{q+1})^{-1}e_{q+1}$:
$r(a)<0<r(c)$.  Its characteristic polynomial
$(t-d)\pi_{q+1}(t)-\gamma\pi_q(t)$ vanishes at both endpoints.
Orthogonality to $\mathcal P_{q-1}$ identifies it with
$(t-a)(t-c)\phi_q^{ac}(t)$.

For an independent matrix proof of exactness expand
$e_1^\top J^ke_1$ as a finite sum over index walks.
A walk starting at index $0$ cannot visit index $q$ and return
using fewer than $2q$ edges; a visit to its diagonal requires at
least $2q+1$ steps.  Hence the changed last diagonal in (AQ7)
preserves moments through degree $2q$.
Similarly a walk using the new coupling of (AQ8) has length at
least $2q+2$, proving moment matching through degree $2q+1$.
This uses finite matrix powers, not a quadrature convergence theorem.
Spectral diagonalization supplies the positive weights
$L(1)|e_1^\top u_i|^2$.

\subsection*{Remainders and the repaired signs}
Repeat the Rolle argument with a simple interpolation condition at
each prescribed endpoint and double conditions at all interior nodes.
Writing the remainder as integral minus quadrature, obtain
\begin{align}
 R_a(f)&=\frac{f^{(2q+1)}(\xi)}{(2q+1)!}
            \int(t-a)(\phi_q^a(t))^2\,d\mu(t),\tag{AQ9}\\
 R_c(f)&=\frac{f^{(2q+1)}(\xi)}{(2q+1)!}
            \int(t-c)(\phi_q^c(t))^2\,d\mu(t),\tag{AQ10}\\
 R_L(f)&=\frac{f^{(2q+2)}(\xi)}{(2q+2)!}
            \int(t-a)(t-c)(\phi_q^{ac}(t))^2\,d\mu(t).\tag{AQ11}
\end{align}
One can first prove the min/max derivative inequalities; the
single-$\xi$ conclusion follows exactly as in (AQ6), with the
nonnegative weight or its negative.  Zero weighted norm means exactness.
Thus a positive odd derivative makes \emph{left Radau a lower bound}
and \emph{right Radau an upper bound}.  The catalogue states these two
directions in reverse.  A positive even derivative makes Lobatto
an upper bound, as the catalogue correctly states.
For $1/t$ and $1/t^2$ on a positive interval the odd derivatives are
negative, so the useful inverse-form bounds are
$Q_c(f)\le L(f)\le Q_a(f)$, alongside the Gauss lower bound.

The discrete orthogonality is an immediate consequence of exactness:
$Q_L(pr)=L(pr)$ whenever $\deg p+\deg r\le2q+1$.
For the Chebyshev weight $(1-t^2)^{-1/2}dt$, put $N=q+1$;
the Lobatto nodes are $\cos(j\pi/N)$, $0\le j\le N$,
with endpoint weights $\pi/(2N)$ and other weights $\pi/N$.
Exactness through degree $2N-1$ follows by summing
$\cos(kj\pi/N)$ with half endpoint weights: the sum is zero
for $1\le k\le2N-1$, while the constant sum is $N$.
This finite geometric-sum verification also covers $N=1$.

\section{Christoffel kernel polynomials and the finite boundary correction}
\label{sec:christoffel-kernel}
\textbf{Atlas item:} \code{kernel-polynomials-christoffel}.
Let $L$ be a positive or quasi-definite polynomial functional,
with monic orthogonal polynomials $\pi_j$ and nonzero squared norms
$h_j=L(\pi_j^2)$.  For a complex shift $z$, use the bilinear
complexification of $L$, not a conjugate-linear inner product.
Suppose $\pi_j(z)\ne0$ through every index actually used, including
the one succeeding the highest requested degree.  Put
\begin{equation}
 q_j=-\frac{\pi_{j+1}(z)}{\pi_j(z)},\quad
 \widehat\pi_j(t)=\frac{\pi_{j+1}(t)+q_j\pi_j(t)}{t-z},\quad
 \widehat L(p)=L((t-z)p).
 \tag{AQ12}
\end{equation}
The numerator vanishes at $t=z$, so polynomial division yields
a monic polynomial of degree $j$.  For $\deg r<j$,
$\widehat L(\widehat\pi_jr)=0$, and
\begin{equation}
 \widehat h_j=\widehat L(\widehat\pi_j^2)=q_jh_j\ne0.
 \tag{AQ13}
\end{equation}
Indeed the $\pi_{j+1}$ term is orthogonal to $\widehat\pi_j$, and
$L(\pi_j\widehat\pi_j)=h_j$.  This proves existence and
quasi-definiteness to the stated order by finite algebra.

Set $e_{-1}=0$ and $e_j=\beta_{j+1}/q_j$.  Evaluating the original
recurrence at $z$ gives
$\alpha_j=z+q_j+e_{j-1}$.  Comparing the next-to-leading coefficients
of the monic polynomials in (AQ12), and using (AQ13), gives
\begin{equation}
 \widehat\alpha_j=z+q_j+e_j,
 \qquad \widehat\beta_j=q_je_{j-1}\quad(j\ge1).
 \tag{AQ14}
\end{equation}
For the first identity, if $s_j$ is the coefficient of $t^{j-1}$
in $\pi_j$, then the corresponding coefficient in
$\widehat\pi_j$ is $s_{j+1}+q_j+z$; subtract consecutive
coefficients and use $\alpha_{j+1}=z+q_{j+1}+e_j$.

Here the nonsymmetric monic Jacobi matrix has diagonal $\alpha_j$,
superdiagonal $1$ and subdiagonal $\beta_j$.
With $L_n$ unit lower bidiagonal of subdiagonal $e_{j-1}$ and
$U_n$ upper bidiagonal of diagonal $q_j$ and superdiagonal $1$,
\begin{equation}
 J_n^{\rm mon}-zI=L_nU_n,\qquad
 \widehat J_n^{\rm mon}=U_nL_n+zI+
                 e_{n-1}E_{n-1,n-1}\quad(n\ge1).
 \tag{AQ15}
\end{equation}
The last rank-one term is essential: finite multiplication drops
the path through index $n$, so the last diagonal of $U_nL_n$
is $q_{n-1}$ rather than $q_{n-1}+e_{n-1}$.
The uncorrected identity is valid for the infinite formal bidiagonal
arrays, or for leading blocks computed from a one-index larger
factorization.  Symmetric Jacobi matrices require the usual diagonal
similarity using the norms $h_j$; the displayed $LU$ factors refer
to the monic matrix.  At dimension zero the identities are empty.

\textbf{Nonvanishing correction.}  The assertion
$\widehat\pi_n(z)\ne0$ is false for general complex $z$, even for
a positive real measure.  For the standard Gaussian probability
measure, $\pi_1=t$, $\pi_2=t^2-1$, and $z=i$ give
$\widehat\pi_1(t)=t-i$, so $\widehat\pi_1(i)=0$, while every
$\pi_k(i)\ne0$ because the original roots are real.
For real $z$ strictly outside the support interval, $(t-z)\mu$
has constant sign; after an overall sign change it is positive.
Its orthogonal-polynomial zeros are in the support interval, and
the additional nonvanishing is then valid.

For SPD $A$ and $z=0$, the transformed functional is
$\widehat L(p)=r_0^\top A p(A)r_0$.  The vectors
$\widehat\pi_j(A)r_0$ are mutually $A$-orthogonal and belong to
successive Krylov spaces with nonzero leading coefficients.
Consequently they are the CG search directions up to nonzero
individual normalizations, by uniqueness of orthogonalization in
each graded space.  This is the appropriate mathematical statement;
the actual CG normalization must be supplied separately.
The symmetric factorization has pivots
$d_0=\alpha_0$, $d_j=\alpha_j-a_j^2/d_{j-1}$ and
subdiagonal $a_j/d_{j-1}$, all well defined before termination by
positive definiteness.  Existing CG span and conjugacy results can
then identify the coefficients without a new algorithm proof.

\section{Second-kind polynomials and Markov convergence without continued-fraction machinery}
\label{sec:second-kind-markov}
\textbf{Atlas item:} \code{jacobi-fraction-second-kind}.
Let $\mu$ be finite and positive, supported in $[a,c]$, with
$\mu_0=L(1)>0$.  For every degree for which $\pi_n$ is defined put
\begin{equation}
 F(z)=\int\frac{d\mu(t)}{z-t},\qquad
 \sigma_n(z)=\int\frac{\pi_n(z)-\pi_n(t)}{z-t}\,d\mu(t),\qquad
 \rho_n(z)=\int\frac{\pi_n(t)}{z-t}\,d\mu(t).
 \tag{AQ16}
\end{equation}
The divided difference in the middle is a polynomial in $(z,t)$,
including at $z=t$; hence $\sigma_n$ is genuinely a polynomial of
degree at most $n-1$.  Direct substitution of the monic recurrence gives
\begin{equation}
 \sigma_0=0,\quad \sigma_1=\mu_0,\quad
 \sigma_{n+1}=(z-\alpha_n)\sigma_n-\beta_n\sigma_{n-1}
 \quad(n\ge1).
 \tag{AQ17}
\end{equation}
The extra term is $L(\pi_n)$, which vanishes for $n\ge1$.
Thus $\sigma_n=\mu_0\pi_{n-1}^{[1]}$, where the associated monic
polynomials use the coefficients shifted by one index.
This establishes the numerator identity by a finite recurrence.

If $\tau_i$ are the roots of $\pi_n$, then
\begin{equation}
 \sigma_n(\tau_i)=\pi_n'(\tau_i)L(\ell_i),\qquad
 \frac{\sigma_n(z)}{\pi_n(z)}
       =\sum_{i=1}^n\frac{\lambda_i}{z-\tau_i}.
 \tag{AQ18}
\end{equation}
The first formula follows by cancellation inside the divided difference.
For the second, multiply by $\pi_n$: two polynomials of degree at most
$n-1$ agree at all $n$ roots.  This proof avoids a general partial-fraction
theorem.  For $n\ge1$, equivalently, a cofactor expansion of the finite Jacobi
resolvent gives
$\sigma_n/\pi_n=\mu_0e_1^\top(zI-J_n)^{-1}e_1$.
Its spectral expansion is valid at every complex $z$ outside the
spectrum of $J_n$.

\textbf{Division correction.}  A literal nested continued fraction
requires the inner denominators to be nonzero, whereas the polynomial
ratio and the full resolvent need only their final denominator nonzero.
For example, $J_2=\left(\begin{smallmatrix}1&1\\1&0\end{smallmatrix}\right)$
is invertible at $z=0$, and the resolvent $(0I-J_2)^{-1}_{11}$ is zero,
but its displayed inner fraction contains $1/0$.
Formalize continued fractions as numerator/denominator pairs satisfying
(AQ17), or add the tail-resolvent nonvanishing hypotheses.
They hold automatically for nonreal $z$ and for real $z$ outside
the full support interval.  The existing real theorem
\code{NLAlib.inv_smul_one_sub_jacobiMatrix_apply_zero_zero}
uses a positive-definiteness hypothesis precisely to avoid these poles.

Two exact identities are particularly useful:
\begin{equation}
 F(z)-\frac{\sigma_n(z)}{\pi_n(z)}
  =\frac{\rho_n(z)}{\pi_n(z)}
  =\frac{1}{\pi_n(z)^2}\int\frac{\pi_n(t)^2}{z-t}\,d\mu(t),
 \qquad z\notin[a,c].
 \tag{AQ19}
\end{equation}
The second follows because
$(\pi_n(z)-\pi_n(t))/(z-t)$ has degree $n-1$ in $t$, so its
product with $\pi_n(t)$ integrates to zero.  These are finite polynomial
orthogonality identities, even though $F$ itself is a transform.

Moment matching and a finite geometric identity also give an explicit
expansion at infinity.  Let $R\ge\max(|a|,|c|)$, $|z|>R$.
Subtract the first $2n$ terms of $1/(z-t)$ from each of the two
equal-moment measures in (AQ18); the remainder is
$t^{2n}/(z^{2n}(z-t))$.  Therefore
\begin{equation}
 \left|F(z)-\frac{\sigma_n(z)}{\pi_n(z)}\right|
 \le \frac{2\mu_0 R^{2n}}{|z|^{2n}(|z|-R)}.
 \tag{AQ20}
\end{equation}
This proves the stated $O(z^{-2n-1})$ with an explicit bound.

\subsection*{A finite Bernstein proof of locally uniform Markov convergence}
Assume $a<c$ and let $K\subset\C\setminus[a,c]$ be nonempty and compact. Put
$\delta=\operatorname{dist}(K,[a,c])>0$.
For $z\in K$, the function $g_z(u)=1/(z-a-(c-a)u)$ on $[0,1]$
has Lipschitz constant at most $(c-a)/\delta^2$.
For an integer $N\ge1$ its Bernstein polynomial is
\[
 B_Ng_z(u)=\sum_{j=0}^N g_z(j/N)\binom Nj u^j(1-u)^{N-j}.
\]
The binomial weights sum to one, and their variance for $j/N$
is $u(1-u)/N$.  Finite Cauchy--Schwarz therefore gives
\[
 |B_Ng_z(u)-g_z(u)|
 \le\frac{c-a}{\delta^2}\sqrt{\frac{u(1-u)}N}
 \le\frac{c-a}{2\delta^2\sqrt N}.
\]
For $n\ge1$ use $N=2n-1$ and Gauss exactness.  Both measures have mass $\mu_0$,
so
\begin{equation}
 \sup_{z\in K}\left|F(z)-\frac{\sigma_n(z)}{\pi_n(z)}\right|
 \le\frac{\mu_0(c-a)}{\delta^2\sqrt{2n-1}}\longrightarrow0.
 \tag{AQ21}
\end{equation}
This proves the full locally uniform conclusion, without a theorem
about infinite continued fractions or complex polynomial approximation.
An empty compact set gives a vacuous uniform assertion; for $a=c$
the measure has one support point and its one-point rule is exact.
Infinite support is required if the sequence is to continue in all
degrees.  At finite support size $N_0$, the $N_0$-point rule is already
exact, and the rational approximants can be extended by the constant
value $F$ thereafter.  Higher orthogonal polynomials with positive norm
must not be invented after this terminal index.

\section{Continuity from a fixed finite moment matrix}
\label{sec:recurrence-continuity}
\textbf{Atlas item:} \code{recurrence-coefficient-continuity}.
Fix $n\ge1$.  In fact only moments through degree $2n-1$ are needed
for the $n$-point rule and its Jacobi coefficients.
Let $m_j^{(N)}\to m_j$ for these finitely many indices, and assume
the limiting moment matrix
$M_n=(m_{i+j})_{0\le i,j<n}$ is positive definite.
Positive definiteness is open: if its smallest eigenvalue is $\eta>0$,
then $\|M_n^{(N)}-M_n\|<\eta/2$ implies
$M_n^{(N)}\succeq(\eta/2)I$.
Thus all relevant polynomials and norms exist for all sufficiently
large $N$, even if some early measures have too few support points.

The coefficient vector of
$\pi_k=t^k+\sum_{j<k}c_jt^j$, $1\le k\le n$, is
\begin{equation}
 c=-M_k^{-1}(m_k,\ldots,m_{2k-1})^\top.
 \tag{AQ22}
\end{equation}
Every entry is a rational function of finitely many moments whose
denominator stays nonzero.  Hence these coefficients converge.
The formulas
$h_j=L(\pi_j^2)$,
$\alpha_j=L(t\pi_j^2)/h_j$ and
$\beta_j=h_j/h_{j-1}$ then give coefficient convergence for $j<n$;
include $\beta_0=\mu_0$ separately if that convention is used.
For $R=\max(1,|a|,|c|)$,
\[
 \sup_{t\in[a,c]}|\pi_n^{(N)}(t)-\pi_n(t)|
 \le\sum_{j=0}^{n-1}|c_j^{(N)}-c_j|R^j\longrightarrow0.
\]

The corresponding symmetric Jacobi matrices converge in norm,
since square root is continuous on their positive subdiagonal squares.
Ordered eigenvalue perturbation bounds imply convergence of the ordered
Gauss nodes.  Simplicity of the limiting nodes means their gaps are
eventually bounded away from zero.  Recover the weights from the
Vandermonde equations
\begin{equation}
 \sum_{i=1}^n\lambda_i\tau_i^k=m_k,
 \qquad 0\le k<n.
 \tag{AQ23}
\end{equation}
The Vandermonde determinant is a product of nonzero node differences,
so inversion is continuous and the weights converge.
This avoids choosing eigenvectors continuously or defining spectral
projectors along a sequence.
At $n=0$ all coefficient, node and weight assertions are empty.
\textbf{Mathematical reduction:} finite matrix inversion, positivity
of a limiting Gram matrix, ordered symmetric eigenvalue continuity,
and one Vandermonde solve.  Weak convergence of measures is unnecessary
once the specified moments converge.

\section{Contour kernels as the difference of two resolvents}
\label{sec:gauss-contour-kernel}
\textbf{Atlas item:} \code{gauss-remainder-contour-kernel}.
Let $\Gamma$ be a positively oriented piecewise $C^1$ Jordan contour
whose interior contains $[a,c]$, and assume $f$ is holomorphic on
an open set containing the contour \emph{and its closed interior}.
The alternative standard hypothesis is that $\Gamma$ is a cycle
null-homologous in the holomorphy domain and has winding number one
at every point of $[a,c]$.
Holomorphy merely near the contour, as in the catalogue hypothesis,
is insufficient for Cauchy's formula.

Cauchy's scalar formula, integrated against $\mu$ and against its
Gauss atomic measure, gives
\begin{equation}
 L(f)-Q_n(f)=\frac1{2\pi i}\oint_\Gamma
 f(z)\left(F(z)-\sum_i\frac{\lambda_i}{z-\tau_i}\right)\,dz
 =\frac1{2\pi i}\oint_\Gamma f(z)\frac{\rho_n(z)}{\pi_n(z)}\,dz.
 \tag{AQ24}
\end{equation}
For a matrix spectral measure this step is a finite sum of scalar
Cauchy formulas, so it needs no interchange theorem for two arbitrary
integrals.  For a general finite measure, put
$\delta_\Gamma=\operatorname{dist}(\Gamma,[a,c])>0$;
the integrand is bounded by
$\max_\Gamma|f|/\delta_\Gamma$ times the contour-speed function.
Its integral is at most
$\mu_0\ell(\Gamma)\max_\Gamma|f|/\delta_\Gamma<\infty$,
which justifies Fubini directly.

Writing $K_n=\rho_n/\pi_n$, the triangle inequality yields exactly
\begin{equation}
 |L(f)-Q_n(f)|\le\frac{\ell(\Gamma)}{2\pi}
                    \max_\Gamma|K_n|\max_\Gamma|f|.
 \tag{AQ25}
\end{equation}
The positive-norm identity (AQ19) additionally yields
$|K_n(z)|\le h_n/(\operatorname{dist}(z,[a,c])|\pi_n(z)|^2)$.
The circle specialization is the first useful formal statement:
\code{NLAlib.eval_interpolate_eq_circleIntegral} and
\code{NLAlib.sub_eval_interpolate_eq_circleIntegral} already expose
Cauchy formulas using \code{circleIntegral};
\code{NLAlib.nodal_resolvent_eq_sum_basis_div} supplies their exact
finite rational identity.  General contours require the corresponding
contour API, but no new quadrature argument.

\section{Moment sensitivity by differentiating orthogonality once}
\label{sec:moment-sensitivity}
\textbf{Atlas item:} \code{recurrence-coefficient-moment-sensitivity}.
Fix $n\ge1$. Let $b_0,\ldots,b_{2n-1}$ be a fixed basis of
$\mathcal P_{2n-1}$, independent of the moments, and let
$m_k=L(b_k)$.  Denote the coefficient of $b_k$ in a polynomial $p$
by $c_k[p]$.  For a perturbation vector $v$ define the linear functional
$\dot L(p)=\sum_k v_kc_k[p]$.
Use the open domain where the Gram matrix on $\mathcal P_{n-1}$
is positive definite.  Formula (AQ22), after a fixed basis change,
makes all quantities below differentiable rational functions of $m$.
No family of representing measures for every nearby moment vector
is necessary.

The product rule for a moment-dependent polynomial is simply
\begin{equation}
 \frac{d}{d\varepsilon}\bigg|_{0} L_\varepsilon(p_\varepsilon)
       =\dot L(p)+L(\dot p).
 \tag{AQ26}
\end{equation}
Since $\pi_j$ is monic, $\deg\dot\pi_j<j$.
Differentiating $L(\pi_j\pi_i)=0$, $i<j$, gives
$L(\dot\pi_j\pi_i)=-\dot L(\pi_j\pi_i)$, because
$L(\pi_j\dot\pi_i)=0$.  Consequently
\begin{equation}
 \dot\pi_j=-\sum_{i<j}\frac{\dot L(\pi_j\pi_i)}{h_i}\pi_i
 \quad(0\le j\le n),\qquad
 \dot h_j=\dot L(\pi_j^2)\quad(0\le j<n).
 \tag{AQ27}
\end{equation}
Differentiate $\beta_j=h_j/h_{j-1}$ to obtain, for $1\le j<n$,
\begin{equation}
 \dot\beta_j=\dot L(\Psi_{2j}),\qquad
 \Psi_{2j}=\beta_j\left(\frac{\pi_j^2}{h_j}
                     -\frac{\pi_{j-1}^2}{h_{j-1}}\right).
 \tag{AQ28}
\end{equation}
For $0\le j<n$ and $\alpha_j=L(t\pi_j^2)/h_j$,
the quotient rule and (AQ27) give
\begin{align}
 \dot\alpha_j&=\dot L(\Psi_{2j+1}),\tag{AQ29}\\
 \Psi_{2j+1}
 &=\frac{(t-\alpha_j)\pi_j^2-2\beta_j\pi_{j-1}\pi_j}{h_j}
  =\frac{\pi_j\pi_{j+1}}{h_j}
       -\frac{\pi_{j-1}\pi_j}{h_{j-1}}.\tag{AQ30}
\end{align}
The terms indexed $j-1$ are omitted at $j=0$.
To verify the only cross term, use
$t\pi_j=\pi_{j+1}+\alpha_j\pi_j+\beta_j\pi_{j-1}$:
$L(t\pi_j\dot\pi_j)=-\beta_j\dot L(\pi_j\pi_{j-1})$.
Thus the entire sensitivity result is a differentiation of finite
linear systems and their orthogonality identities.

In orthonormal notation $p_j=\pi_j/\sqrt{h_j}$ these become
$\Psi_{2j}=\beta_j(p_j^2-p_{j-1}^2)$ and
$\Psi_{2j+1}=a_{j+1}p_jp_{j+1}-a_jp_{j-1}p_j$,
when those normalized polynomials exist.
\textbf{Indexing repairs:} $\beta_0=L(1)$ has
$\dot\beta_0=\dot L(1)$; (AQ28) is not its formula.
For the highest requested diagonal coefficient $\alpha_{n-1}$,
use the monic expression (AQ30).  It requires only moments through
$2n-1$; the expression involving $p_n$ and $a_n$ separately would
apparently require $h_n$ and the additional moment $m_{2n}$,
although their product cancels that dependence.

Taking $v$ to be the $k$-th coordinate vector proves
\begin{equation}
 \frac{\partial\alpha_j}{\partial m_k}=c_k[\Psi_{2j+1}],\qquad
 \frac{\partial\beta_j}{\partial m_k}=c_k[\Psi_{2j}]
 \quad(1\le j<n),\qquad
 \frac{\partial\beta_0}{\partial m_k}=c_k[1].
 \tag{AQ31}
\end{equation}
Here the $\alpha_j$ identity is for $0\le j<n$ and every moment
coordinate has $0\le k<2n$.
After putting rows in the declared output order, the induced norm
of this Jacobian is the \emph{absolute} local condition number.
The relative condition number also has the factor
$\|m\|/\|\Phi(m)\|$ whenever the latter normalization is defined.
The verified bridge
\code{NLAlib.specNorm_euclideanJacobian} already identifies the
Euclidean Jacobian norm with \code{NLAlib.mapAbsCondition}.

\section{Rational measure modification as a finite constrained polynomial space}
\label{sec:christoffel-uvarov}
\textbf{Atlas item:} \code{christoffel-uvarov-modification}.
Let $u,v\in\R[t]$ be coprime, of degrees $\ell,m$, with simple
complex roots $u_1,\ldots,u_\ell$ and $v_1,\ldots,v_m$, no roots
on the compact support interval, and $u/v>0$ there.
Suppose the positive modified functional
$\widehat L(p)=L(up/v)$ is positive definite on $\mathcal P_{n-1}$.
This suffices for a unique degree-$n$ monic $\widehat\pi_n$.
The following finite system gives a direct route to the determinant
formula, including $m>n$, without a chain of modification algorithms.

Set $P=u\widehat\pi_n$.  It is characterized by
\begin{equation}
 \deg P\le n+\ell,\qquad
 P(u_i)=0\ (1\le i\le\ell),\qquad
 L(t^rP/v)=0\ (0\le r<n),\qquad
 [t^{n+\ell}]P=\operatorname{lc}(u).
 \tag{AQ32}
\end{equation}
Necessity is immediate.  Conversely the root conditions imply
$u\mid P$, so $P/u$ is monic of degree $n$ and satisfies exactly
the defining orthogonality for $\widehat L$.
The homogeneous system obtained by replacing the last condition by
zero has only the zero solution: then $P=u s$ with $\deg s<n$,
and its equations give $\widehat L(s^2)=0$, hence $s=0$.
For complex coefficients apply the same argument to
$\widehat L(s\overline s)>0$; the equations annihilate every
complex polynomial of degree $<n$, including $\overline s$.
Therefore the square coefficient matrix is invertible.  This is
the exact reason that its normalization determinant does not vanish.

Take any fixed polynomial basis $p_0,\ldots,p_{n+\ell}$ with
$\deg p_j=j$ and leading coefficient one (monomials always work).
Form the $(n+\ell+1)\times(n+\ell+1)$ determinant $D(t)$ whose
first row is $(p_0(t),\ldots,p_{n+\ell}(t))$, whose next $\ell$
rows are evaluations at the roots $u_i$, and whose remaining $n$
rows are
\[
 \left(L(t^r p_j(t)/v(t))\right)_{0\le j\le n+\ell},
 \qquad 0\le r<n.
\]
Let $d=[t^{n+\ell}]D(t)$.  The preceding invertibility argument
proves $d\ne0$ and
\begin{equation}
 u(t)\widehat\pi_n(t)=\operatorname{lc}(u)\,\frac{D(t)}d.
 \tag{AQ33}
\end{equation}
This bordered formula works for every $m$, including $m>n$, and
uses no negative-index orthogonal polynomial.
Choosing monomials also avoids requesting original orthogonal
polynomials past a finite-support termination index.

When $m\le n$ and the original polynomials up to $n+\ell$ exist,
there is a smaller familiar formula.  Orthogonality applied to
$vr$, $\deg r<n-m$, shows $L(Pr)=0$; hence
\[
 P\in\operatorname{span}\{\pi_{n-m},\ldots,\pi_{n+\ell}\}.
\]
For $0\le r<n$, divide $t^r$ by $v$ and use partial fractions of
the remainder.  The remaining conditions in (AQ32) reduce exactly to
\begin{equation}
 P(u_i)=0,\qquad
 L\!\left(\frac{P(t)}{v_j-t}\right)=0
 \quad(1\le j\le m).
 \tag{AQ34}
\end{equation}
Indeed the polynomials $v(t)/(t-v_j)$ form a basis of
$\mathcal P_{m-1}$, as evaluation at each distinct root proves.
Thus (AQ33) reduces to a determinant with first row
$\pi_{n-m},\ldots,\pi_{n+\ell}$, evaluation rows at $u_i$,
and Cauchy-transform rows $\rho_k(v_j)$, exactly the catalogue form.

For $m>n$, all the functionals in the larger bordered formula can
still be computed from Cauchy transforms.  For the chosen basis put
$\eta_k(z)=L(p_k(t)/(z-t))$; these are $\rho_k$ when $p_k=\pi_k$.
Since $r<m$,
\begin{equation}
 \frac{t^r}{v(t)}=\sum_{j=1}^m
       \frac{v_j^r}{v'(v_j)(t-v_j)},\qquad
 L\!\left(\frac{t^rp_k(t)}{v(t)}\right)
   =-\sum_{j=1}^m\frac{v_j^r\eta_k(v_j)}{v'(v_j)}.
 \tag{AQ35}
\end{equation}
Multiply the rational identity by $v$ and test its values at its
$m$ roots to prove it.  Complex conjugate roots cause no difficulty:
the unique normalized solution is real because conjugation preserves
all the real coefficient equations.

For a quadratic divisor, $u=1$, $m=2$, and $n\ge2$, (AQ34) yields
\[
 \widehat\pi_n=\pi_n+s_n\pi_{n-1}+t_n\pi_{n-2},\qquad
 \begin{pmatrix}\rho_{n-1}(v_1)&\rho_{n-2}(v_1)\\
                 \rho_{n-1}(v_2)&\rho_{n-2}(v_2)\end{pmatrix}
 \binom{s_n}{t_n}=-\binom{\rho_n(v_1)}{\rho_n(v_2)}.
\]
Its determinant is nonzero by the same uniqueness argument.
For $n<2$ use (AQ32), rather than nonexistent negative indices.
Once $\widehat\pi_j$ are known, the modified recurrence coefficients
are simply $\widehat h_j/\widehat h_{j-1}$ and
$\widehat L(t\widehat\pi_j^2)/\widehat h_j$.
The derivation establishes these exact identities, not the numerical
stability of any update algorithm.

\section{Rational Gauss quadrature by clearing a fixed denominator}
\label{sec:rational-gauss}
\textbf{Atlas item:} \code{rational-gauss-quadrature}.
Let $n\ge1$ and
$\omega(t)=\prod_{j=1}^{r}(1+\zeta_jt)^{s_j}$ have degree
$m=\sum_js_j\le2n$, with distinct nonzero $\zeta_j$ and no zeros
on $[a,c]$.  Repeated equal $\zeta_j$ are combined first.
Suppose $d\nu=d\mu/\omega$ is a positive measure and its $n$-point
Gauss rule is $(\tau_i,\lambda_i^G)$.
Define $w_i=\omega(\tau_i)\lambda_i^G$.
For every polynomial $p$ of degree at most $2n-1$,
\begin{equation}
 \int\frac{p}{\omega}\,d\mu
 =\int p\,d\nu
 =\sum_i\lambda_i^G p(\tau_i)
 =\sum_iw_i\frac{p(\tau_i)}{\omega(\tau_i)}.
 \tag{AQ36}
\end{equation}
Conversely, exactness on this rational space implies polynomial
exactness for the weights $w_i/\omega(\tau_i)$ and measure $\nu$,
so the uniqueness theorem identifies that rule as Gauss.
This is a change of measure and a finite polynomial identity.

Polynomial division and partial fractions prove the exact space equality
\begin{equation}
 \left\{\frac p\omega:\deg p\le2n-1\right\}
 =\mathcal P_{2n-1-m}\ \oplus\
   \operatorname{span}\{(1+\zeta_jt)^{-s}:1\le s\le s_j\},
 \tag{AQ37}
\end{equation}
with $\mathcal P_{-1}=\{0\}$ when $m=2n$.
The directness follows by clearing denominators and examining the
principal part at each pole; alternatively successively evaluate
the cleared numerator and its derivatives at that pole.
For $m=0$ this is ordinary polynomial Gauss quadrature.
For $n=0$ necessarily $m=0$, the space is $\{0\}$ and the assertion
is vacuous.

\textbf{Positivity correction.}  Nonvanishing of $\omega$ alone
does not make $\mu/\omega$ a positive measure when complex poles
are permitted.  The algebraic equivalence (AQ36) still makes sense
for a complex functional if a polynomial-exact rule is supplied,
but positive Gauss weights and real interior nodes need positivity
as an additional hypothesis.  A real nonvanishing denominator has
constant sign on the interval; one may divide by $|\omega|$ and
restore that constant sign.  A conjugate-paired denominator positive
on the interval is another sufficient hypothesis.  None of these
claims requires constructing or analyzing a rational Krylov algorithm.


\section{Crouzeix--Palencia: finite matrix algebra around one analytic interface}
\label{math:cp}
\paragraph{Catalogue item.} \code{crouzeix-palencia}.
The target is
\[
 \|p(A)\|_2\le(1+\sqrt2)\max_{z\in W(A)}|p(z)|.
\]
For dimension zero the left side is zero; state that case separately
instead of taking a maximum over the empty numerical range.
For positive dimension the following isolates precisely the part that
cannot presently be replaced by finite algebra.

\begin{lemma}[An algebra homomorphism estimate]
Let $\mathcal B$ be a unital normed algebra of scalar functions with
the supremum norm, and let $\Phi:\mathcal B\to\C^{d\times d}$ be a
bounded unital algebra homomorphism, with $d\ge1$. Suppose that for every $f$ there
is $g\in\mathcal B$ with
\[
 \|g\|\le\|f\|,\qquad
 \|\Phi(f)+\Phi(g)^*\|_2\le2\|f\|.
                                                        \tag{CP1}
\]
Then $\|\Phi\|\le1+\sqrt2$.
\end{lemma}
\begin{proof}
Put $K=\|\Phi\|\ge1$, and for $\|f\|\le1$ write
$F=\Phi(f)$, $G=\Phi(g)$, $S=F+G^*$. Multiplicativity gives
\[
 (FF^*)^2=FS^*FF^*-\Phi(fgf)F^*.
\]
Consequently $\|F\|^4\le2K^3+K^2$. Taking the supremum over the
unit ball yields $K^4\le2K^3+K^2$, so $K^2\le2K+1$.
\end{proof}
This is the short algebra argument of
\href{https://arxiv.org/abs/1708.08633}{Ransford--Schwenninger};
it requires no extremizing function or compactness of a function ball.

\paragraph{A finite definition of the evaluation homomorphism.}
Factor the minimal polynomial over $\C$ as
$m_A(z)=\prod_{j=1}^s(z-\lambda_j)^{r_j}$.
For a function holomorphic near the spectrum, let $H_f$ be the unique
polynomial of degree less than $\deg m_A$ satisfying
\[
 H_f^{(\ell)}(\lambda_j)=f^{(\ell)}(\lambda_j)
 \quad(0\le\ell<r_j).
\]
Confluent interpolation, or polynomial Chinese remainders, gives
existence and uniqueness. Define $\Phi(f)=H_f(A)$.
Leibniz's finite rule shows $H_{fg}\equiv H_fH_g\pmod {m_A}$;
therefore $\Phi$ is a unital homomorphism and agrees with polynomial
evaluation. This definition uses a finite list of scalar derivatives,
not a general functional calculus for nonnormal operators.
If the spectrum lies inside a bounded domain $\Omega$, choose a small
closed disk around each $\lambda_j$ contained in $\Omega$.
Scalar Cauchy estimates bound every required jet by a constant times
$\|f\|_{\overline\Omega}$. The fixed confluent interpolation matrix
then proves boundedness of $\Phi$.

\paragraph{The indispensable analytic lemma, stated explicitly.}
For a smooth bounded convex domain $\Omega$ containing $W(A)$ in its
interior, let $f$ be holomorphic on $\Omega$ and continuous on
$\overline\Omega$. The scalar transform
\[
 g(z)=\frac1{2\pi i}\int_{\partial\Omega}
             \frac{\overline{f(\zeta)}}{\zeta-z}\,d\zeta
\]
extends continuously to $\overline\Omega$ and satisfies
$\|g\|_{\overline\Omega}\le\|f\|_{\overline\Omega}$.
The matrix density, with $\nu$ the outward unit normal and
$R=(\zeta I-A)^{-1}$, is
\[
 M(\zeta)=\frac{\nu R+\overline\nu R^*}{2\pi}\succeq0,
 \qquad \int_{\partial\Omega}M\,ds=2I,
 \qquad \Phi(f)+\Phi(g)^*=\int_{\partial\Omega}fM\,ds.
                                                        \tag{CP2}
\]
The positivity itself is finite algebra:
\[
 2\pi M
 =R^*\{\nu(\overline\zeta I-A^*)+
                 \overline\nu(\zeta I-A)\}R\succeq0,
\]
by the supporting half-plane inequality for $W(A)$.
Matrix Cauchy--Schwarz and $\int M=2I$ imply
$\|\int fM\|\le2\|f\|$, giving (CP1).
The scalar boundary-extension/contraction theorem and Cauchy identities
in (CP2) are the analytic content of
\href{https://arxiv.org/abs/1702.00668}{Crouzeix--Palencia}.
Shrinking such convex neighborhoods to the compact set $W(A)$ gives
the polynomial target, by uniform continuity of $p$.

\paragraph{Audit status.}
The homomorphism lemma and jet construction above are complete
reductions. The scalar boundary theorem, its matrix Cauchy bridge,
and the smooth-neighborhood approximation still require analytic
formalization. Merely postulating (CP1) for polynomial evaluation
would hide the main difficulty: the partner $g$ need not be a
polynomial. No proof using only finite linear algebra is claimed here.

\section{The two sharp approximation entries: reusable subcases and exact residual tasks}
\label{math:sharp}
\paragraph{Catalogue items.}
\code{abs-best-approx-asymptotics}, \code{exp-rational-halfline}.
These entries bundle several different claims:
\[
 nE_{n,0}(|x|;[-1,1])\longrightarrow\beta,
 \quad e^{\pi\sqrt n}E_{n,n}(|x|;[-1,1])\longrightarrow8,
\]
\[
 E_{0,n}(e^{-x};[0,\infty))^{1/n}\longrightarrow\tfrac13,
 \quad \frac{E_{n,n}(e^{-x};[0,\infty))}
 {2H^{n+1/2}}\longrightarrow1.
\]
Here $0<H<1$ is Halphen's constant in the convention
$H\simeq1/9.28902549$. The decimal is not its formal definition.
The earlier document
\code{polynomial_approximation_operator_derivations.tex}
contains a coefficient-compactness proof of the polynomial Bernstein
limit and a finite Laguerre Gram proof of the reciprocal root limit.
Those are mathematical derivations, not declarations found in the
current Lean index. Both arguments are reconstructed below in a shorter
form. The present note does not upgrade their status to machine-checked
results.

\subsection{The polynomial Bernstein limit by subtracting a comparison polynomial}
Write $E_n=E_{n,0}(|x|;[-1,1])$. Symmetrization shows
$E_{2m+1}=E_{2m}$. For even $n=2m$, the known Chebyshev expansion gives
\[
 S_n(x)=\frac2\pi+\frac4\pi\sum_{j=1}^m
       \frac{(-1)^{j+1}T_{2j}(x)}{4j^2-1},\qquad
 \|S_n-|x|\|=\frac2{\pi(n+1)}.              \tag{B1}
\]
The upper estimate follows by summing the absolute coefficient tail;
at $x=0$ the omitted terms have the same sign, proving equality.

A finite dual certificate gives positivity of the eventual constant.
For $K=\lfloor n/2\rfloor+1$ set
\[
 F_K(t)=\frac1{K+1}\left|\sum_{j=0}^K e^{ijt}\right|^2,
 \qquad
 \Lambda(h)=\frac1{2\pi}\int_{-\pi}^{\pi}
             h(\sin t)\cos(4Kt)F_K(2t)\,dt.
\]
$F_K\ge0$ has normalized integral one, so $\|\Lambda\|\le1$.
The multiplier has only frequencies with absolute value at least
$2K>n$, whereas $p(\sin t)$ has frequencies at most $n$ for
$\deg p\le n$. Hence $\Lambda(p)=0$.
Direct integration of $|\sin t|$ gives its nonzero even Fourier
coefficients $-2/(\pi(4j^2-1))$. Writing
$w_j=1-|j|/(K+1)$ gives the finite calculation
\[
 |\Lambda(|x|)|=
 \sum_{j=-K}^K\frac{2w_j}{\pi(4(2K+j)^2-1)}
 \ge\frac{2(K+1)}{\pi(36K^2-1)}
 \ge\frac1{18\pi(n+1)}.                    \tag{B2}
\]
Thus $1/(18\pi(n+1))\le E_n\le2/(\pi n)$ for $n\ge1$.

Put $B_n(z)=nS_n(z/n)$ for even $n$. Let $a$ be the local complex
inverse of sine at zero, so $a(0)=0,a'(0)=1$.
For each fixed disk, $n a(z/n)\to z$ uniformly; this follows directly
from $a(u)-u=o(|u|)$. The finite Chebyshev recurrence gives
$T_{2j}(\sin u)=(-1)^j\cos(2ju)$ and therefore
\[
 B_n(z)=\frac{2n}{\pi(n+1)}+
 \frac{4n}{\pi}\sum_{j=1}^{n/2}
       \frac{1-\cos(2ja(z/n))}{4j^2-1}.     \tag{B3}
\]
Use the continuous extension
\[
 \frac{1-\cos(tz)}{t^2}
 =z^2\int_0^1(1-s)\cos(stz)\,ds
\]
at $t=0$. With mesh $2/n$, (B3) is a Riemann sum whose extra
weights are $4j^2/(4j^2-1)$. Their error is $O(1/n)$ uniformly
on each fixed disk, since $\sum_{j\ge1}(4j^2-1)^{-1}=1/2$.
Uniform continuity on compact rectangles proves
\[
 B_n\longrightarrow B(z)=\frac2\pi+
       \frac2\pi\int_0^1\frac{1-\cos(tz)}{t^2}\,dt,
 \quad |B(z)|\le\frac{2+|z|^2}{\pi}e^{|\operatorname{Im}z|}. \tag{B4}
\]
The cosine series shows directly that $B$ is entire.

Choose even minimizers $p_n$ and an even subsequence with
$nE_n\to b=\liminf_{n\text{ even}}nE_n$.
Set $D_n(z)=np_n(z/n)-B_n(z)$. By (B1),
$\|D_n\|_{[-n,n]}<4/\pi$ and $\deg D_n\le n$.
For a degree-$n$ polynomial $p$ with $\|p\|_{[-1,1]}\le M$,
apply the disk maximum principle to the polynomial
$v^np((v+v^{-1})/2)$, expanded to remove negative powers.
Substitution $v=1/w$ gives
\[
 |p((w+w^{-1})/2)|\le M|w|^n\quad(|w|\ge1).
\]
For $D_n$ this implies, wherever the chosen local inverse $a(z/n)$
is defined (in particular on each fixed disk for all sufficiently
large $n$),
\[
 |D_n(z)|\le\frac4\pi
          \exp\bigl(n|\operatorname{Im}a(z/n)|\bigr).       \tag{B5}
\]
The right side is bounded on each fixed disk, uniformly for large $n$.
For completeness, coefficient compactness here needs only scalar
compactness: on $|z|=R$, finite Fourier orthogonality bounds the
$k$th coefficient by $M_RR^{-k}$. Take a diagonal subsequence of
coefficients. On any smaller disk the geometric bound controls all
tails, proving locally uniform convergence to an entire $D$.
Equation (B5) passes to the limit as
$|D(z)|\le(4/\pi)e^{|\operatorname{Im}z|}$.
Thus $F=B+D$ is entire, real and even, with
\[
 |F(z)|\le\frac{6+|z|^2}{\pi}e^{|\operatorname{Im}z|},\qquad
 |F(x)-|x||\le b\quad(x\in\R).              \tag{B6}
\]

Finally fix $0<\alpha<1$, set $t=(1-\alpha)/2$ and $\rho=e^t$.
Then $c=t(1-e^{-t})>0$ and
$\alpha\sinh t-t\le-c$, since $\alpha\le e^{-2t}$ and
$\sinh t\le te^t$. On the Bernstein ellipse of parameter $\rho$,
the function $G_N(w)=F(\alpha Nw)/(\alpha N)$ has bound
\[
 M_N=\frac{6+(\alpha N\cosh t)^2}{\pi\alpha N}
                   e^{\alpha N\sinh t}.
\]
Degree-$N$ Chebyshev truncation has error at most
$2M_N\rho^{-N}/(\rho-1)$. On $[-1,1]$, (B6) adds error
$b/(\alpha N)$ relative to $|w|$. Hence
\[
 NE_N\le\frac b\alpha+
 \frac{2(6+(\alpha N\cosh t)^2)}{\pi\alpha(e^t-1)}e^{-cN}.
\]
Take the limsup and then $\alpha\uparrow1$. The odd-even identity
and $E_{2m}\to0$ show the full liminf is $b$ as well. Consequently
\[
 nE_n\longrightarrow\beta=b,\qquad
 \frac1{18\pi}\le\beta\le\frac2\pi.
\]
This proves existence and positivity of the sharp limit; it does not
certify printed decimal digits of $\beta$. The formalization uses the
already available ellipse approximation bound, finite Fourier
orthogonality, scalar subsequence compactness, and a local sine inverse.

\paragraph{Verified reuse for this limit.}
The coefficient expansion and tail in (B1) are supplied by
\code{NLAlib.hasSum_chebyshevCoeff_abs_mul_eval_T},
\code{NLAlib.chebyshevCoeff_abs_even}, and
\code{NLAlib.tsum_abs_chebyshevCoeff_abs_tail}.
The finite polynomial growth estimate preceding (B5) already appears as
\code{NLAlib.norm_eval_le_mul_pow_of_mem_filledBernsteinEllipse}.
The final approximation step can directly use
\code{NLAlib.exists_degree_le_abs_sub_eval_le_of_open_bernsteinEllipse},
with its stated constant $2M\rho^{-N}/(\rho-1)$.
Thus the new work is the scaling/compactness bridge and its finite dual
lower bound, rather than either Chebyshev approximation theorem.

\subsection{An exact finite certificate shared by both rational targets}
Write $E_{m,n}(f;I)$ for pole-free rational minimax error on a compact
real interval. Choose a nonempty finite family of distinct nodes
$t_i\in I$, positive weights
$a_i$, signs $s_i\in\{-1,1\}$ and $\eta\ge0$ such that
\begin{align}
 \sum_i a_i s_i t_i^j&=0 &&(0\le j\le m),\tag{RC1}\\
 \sum_i a_i(s_i f(t_i)-\eta)t_i^j&=0 &&(0\le j\le n).\tag{RC2}
\end{align}
Then $E_{m,n}(f;I)\ge\eta$.

\begin{proof}
Every admissible $P/Q$ can be reduced and scaled so $Q>0$ on $I$.
For $e=\|f-P/Q\|_I$, multiply the nodewise inequality
$s_i(f(t_i)Q(t_i)-P(t_i))\le eQ(t_i)$ by $a_i$ and sum.
(RC1)--(RC2) turn its left side into
$\eta\sum_i a_iQ(t_i)$. This sum is positive, so $e\ge\eta$.
Take the infimum. This includes zero numerators and degree defects;
no existence of a nondefective minimizer is assumed.
\end{proof}

For an upper bound it suffices to exhibit $P,Q$, with $Q>0$, and
verify $-\eta Q\le fQ-P\le\eta Q$ on $I$.
For $f(x)=|x|$, these are four polynomial sign inequalities on
$[-1,0]$ and $[0,1]$. These are explicit finite sufficient certificates
for fixed-degree upper and lower bounds. This does not by itself give
certificates uniform in the degree with the sharp asymptotic constant.
The current snapshot already supplies
\code{NLAlib.exists_unique_norm_sub_eq_rationalMinimaxError} and
\code{NLAlib.le_rationalMinimaxError_of_alternation}; use them instead
of reproving minimizer existence or rational alternation.

\subsection{An exact parity reduction for the Stahl clause}
Uniqueness of the best rational approximant to an even function makes
that approximant even. For a reduced rational $P/Q$ with no pole at
zero, evenness implies $P(-x)Q(x)=P(x)Q(-x)$. Coprimality makes
$Q(-x)=cQ(x)$; evaluating at zero gives $c=1$, and then $P$ is even.
Consequently, for every $m\ge0$,
\[
 E_{2m,2m}(|x|;[-1,1])
 =E_{m,m}(\sqrt t;[0,1])
 =E_{2m+1,2m+1}(|x|;[-1,1]).               \tag{RC3}
\]
Both directions follow by the substitution $t=x^2$ and the even
reduced representation. In particular the exact sharp dependency is
the square-root case of
\href{https://arxiv.org/abs/math/9301217}{Stahl's strong asymptotic}:
\[
 e^{\pi\sqrt{2m}} E_{m,m}(\sqrt t;[0,1])\longrightarrow8.
\]
Since $\sqrt{2m+1}-\sqrt{2m}\to0$, it implies the full sequence
in (RC3). The finite parity reduction is complete; the strong
asymptotic is not proved by this reduction.

\subsection{Compactifying the exponential with a quantitative finite surrogate}
Set $t=(1+x)^{-1}$ and
\[
 F(t)=\begin{cases}e^{1-1/t},&t>0,\\0,&t=0.\end{cases}
\]
Substitution and multiplication by $t^n$ carry type-$(n,n)$ rational
functions to type-$(n,n)$ functions. A rational approximant with finite
error on the half-line is bounded at infinity, so its transformed
singularity at $t=0$ is removable. Conversely a pole-free transformed
function is bounded and pole-free on the half-line. Hence
\[
 E_{n,n}(e^{-x};[0,\infty))=E_{n,n}(F;[0,1]). \tag{RC4}
\]
For $0<t\le1$, $F'(t)=t^{-2}e^{1-1/t}\le4/e<2$;
therefore $F$ is $2$-Lipschitz on $[0,1]$.
For an integer $M\ge1$, let its degree-$M$ Bernstein polynomial be
\[
 B_MF(t)=\sum_{j=0}^M F(j/M)\binom Mj t^j(1-t)^{M-j}.
\]
Finite binomial moment identities and Cauchy--Schwarz give
\[
 |B_MF(t)-F(t)|\le2\sqrt{t(1-t)/M}\le M^{-1/2}. \tag{RC5}
\]
For any two target functions on the same domain, the triangle
inequality followed by infima gives
$|E_{m,n}(f)-E_{m,n}(g)|\le\|f-g\|$.
Choosing $M_n=\lceil H^{-4n}\rceil$ in (RC5) gives
\[
 |E_{n,n}(F)-E_{n,n}(B_{M_n}F)|\le H^{2n}=o(H^n). \tag{RC6}
\]
This converts the Halphen equivalent into an equivalent sequence of
finite polynomial-target rational problems. Certified enclosures of
the finitely many coefficients may replace their exact values if the
resulting uniform error is also $o(H^n)$.

To finish by certificates one must construct, for every large $n$,
an upper approximant and a lower certificate with errors
$(2\sqrt H+o(1))H^n$. A lower estimate with an unspecified
exponential base does not suffice. The classical sharp input is
\href{https://doi.org/10.1070/SM2002v193n01ABEH000619}
{Aptekarev's strong-asymptotic theorem}. Neither (RC6) nor the existence
of minimizers supplies those asymptotically sharp certificates.

\subsection{Two shorter subcases that should not remain bundled with that gap}
Polynomial approximation of $e^{-x}$ on $[0,\infty)$ has exact
minimax error $1/2$ at every degree: a nonconstant polynomial has
unbounded error, whereas a constant $c$ has error
$\max\{|c|,|1-c|\}\ge1/2$, attained at $c=1/2$.

For the reciprocal-polynomial clause, the earlier finite Gram route
reduces the entire root limit to
\[
 \frac{3^{-n}}{6\sqrt{(4n+4)\log3+\log4}}
 \le E_{0,n}(e^{-x};[0,\infty))\le\sqrt2\,3^{-n}. \tag{RC7}
\]
Taking logarithms and dividing by $n$ immediately gives $1/3$.
Here is a self-contained reconstruction of (RC7). Scaling $x$ lets
us approximate $e^{-x/4}$ instead. Put $f(x)=e^{x/4}$,
$\langle u,v\rangle=\int_0^\infty uv e^{-x}\,dx$, and
\[
 L_j(x)=\sum_{r=0}^j(-1)^r\binom jr\frac{x^r}{r!}.
\]
The finite Leibniz identity
$L_j=e^x D^j(e^{-x}x^j)/j!$, integrated by parts $j$ times, gives
\[
 \langle L_j,p\rangle=\frac{(-1)^j}{j!}
             \int_0^\infty x^je^{-x}p^{(j)}(x)\,dx.
\]
Boundary terms vanish at zero because the derivatives of $x^je^{-x}$
of order below $j$ vanish, and at infinity by exponential decay.
Using $\int_0^\infty x^re^{-s x}\,dx=r!/s^{r+1}$ gives
\[
 \langle L_i,L_j\rangle=\delta_{ij},\qquad
 c_j=\langle f,L_j\rangle=\frac{4(-1)^j}{3^{j+1}},\qquad
 \langle f,f\rangle=2.
\]
The $L_j$ form a basis of $\mathcal P_n$, by their nonzero triangular
leading coefficients. For $g_n=\sum_{j=0}^nc_jL_j$, finite expansion
of a square therefore proves
\[
 \|f-p\|_{L^2(e^{-x})}^2
 =\frac{2}{3^{2n+2}}+\|p-g_n\|_{L^2(e^{-x})}^2. \tag{RC8}
\]
There is no use of completeness of a Laguerre expansion.

On $\mathcal P_n$, differentiation satisfies $D^{n+1}=0$. Thus
$(I-D)^{-1}=\sum_{j=0}^nD^j$, and integration by parts gives
\[
 ((I-D)^{-1}g)(x)=e^x\int_x^\infty g(t)e^{-t}\,dt.
\]
Set $s_n=\tfrac34(I-D)^{-1}g_n$. Since the same integral expression,
with factor $3/4$, fixes $f$, Cauchy--Schwarz and (RC8) imply
\[
 |f(x)-s_n(x)|\le\frac{\sqrt2}{4\,3^n}e^{x/2}\quad(x\ge0).
                                                               \tag{RC9}
\]
For $a>0$ let $q=q_{n,a}$ minimize
$\max_{[0,a]}e^{-x/2}|f(x)-q(x)|$, and denote the minimum by
$\delta_{n,a}$. Existence follows from bounded values at $n+1$
fixed nodes, invertibility of the Vandermonde matrix, and coefficient
compactness. The usual sign-perturbation proof of alternation works
with any strictly positive continuous weight: if the nonzero
extremal error had at most $n+1$ alternating sign blocks, put roots
between consecutive blocks and perturb by a small polynomial of the
matching sign. Compactness gives a strict improvement, a contradiction.
Thus $f-q$ has $n+1$ distinct zeros $t_0<\cdots<t_n$ in $(0,a)$.

Newton interpolation at these zeros gives
\[
 q(x)=\sum_{j=0}^n f[t_0,\ldots,t_j]
                    \prod_{\ell<j}(x-t_\ell).
\]
Repeated Rolle's theorem gives positivity of every displayed divided
difference, and of the interpolation remainder for $x\ge a$.
Consequently
\[
 0<q(x)\le f(x),\qquad q'(x)\ge0\quad(x\ge a). \tag{RC10}
\]
The derivative assertion includes $n=0$, when it is zero. By (RC9),
$\delta_{n,a}\le\sqrt2/(4\,3^n)$.

For the upper bound choose $a=4n\log3+\log4$. On $[0,a]$,
$\delta_{n,a}e^{x/4}\le1/2$, hence $q\ge f/2$ and
$|f^{-1}-q^{-1}|\le2\delta_{n,a}$.
For $x\ge a$, (RC10) yields
\[
 0\le q(x)^{-1}-f(x)^{-1}\le q(a)^{-1}
 \le2e^{-a/4}=\sqrt2\,3^{-n}.
\]
This denominator is positive on the entire half-line and proves the
upper bound in (RC7).

For the lower bound choose $a=(4n+4)\log3+\log4>2$ and suppose
an admissible reciprocal $1/p$ had error
$\epsilon<1/(2\sqrt a\,3^{n+1})$.
Its value at zero and absence of poles make $p>0$ throughout the
half-line. Since $\epsilon e^{a/4}<1/2$, on $[0,a]$ one has
$p\le2f$ and $|p-f|\le2\epsilon e^{x/2}$.
Optimality of $q_{n,a}$ gives the same last bound for $q$.
Beyond $a$, (RC10) gives $|q-f|\le f$. Therefore (RC8) implies
\[
 \frac2{3^{2n+2}}
 \le\int_0^\infty(q-f)^2e^{-x}\,dx
 \le4a\epsilon^2+2e^{-a/2}
 =4a\epsilon^2+\frac1{3^{2n+2}},
\]
contradicting the chosen $\epsilon$. No half-line minimizer is needed.
Absorbing a nonzero constant numerator into $p$ handles type $(0,n)$;
the zero numerator has error one and cannot beat this lower bound.
This proves (RC7) and the root limit.

\paragraph{Operator consequence.}
For each $n$ the constructed positive denominator $q_n$ satisfies
\[
 \|e^{-A}-q_n(4A)^{-1}\|_2\le\sqrt2\,3^{-n}
\]
for every finite-dimensional PSD $A$, by diagonalization.
The constant is independent of the largest eigenvalue.

\paragraph{What remains open in these proposed routes.}
The outstanding new mathematics is an easier proof of Stahl's
prefactor and the Halphen strong equivalent, or asymptotically sharp
explicit certificates replacing them. The Bernstein-limit and
reciprocal-root subcases have existing mathematical routes to reuse.
They must be tracked separately from the two strong asymptotics.


\appendix
\section{Coverage of the refreshed catalogue}
The following list accounts for every selected entry, including the clauses
that are explicitly retained as dependencies rather than proved here.
\begin{longtable}{@{}p{0.43\textwidth}p{0.53\textwidth}@{}}
\toprule
Catalogue identifier & Route or remaining boundary\\
\midrule\endhead
\code{truncation-lemma} & Projection Pythagoras; unspecified early variant left open\\
\code{xtrace-unbiased} & Disjoint-family independence and an integrable uniform bound\\
\code{positive-real-moment-comparison} & Scalar tangent inequality; existing real-exponent Lp comparison\\
\code{hutchpp-unbiased} & Product-law integration of the trace split\\
\code{krylov-grade} & Annihilator ideal and polynomial division\\
\code{residual-polynomial} & Finite evaluation image and compact sublevel sets\\
\code{field-of-values-def} & Two-dimensional compression and the Bloch sphere\\
\code{crouzeix-palencia} & Finite homomorphism estimate; analytic boundary theorem retained\\
\code{jacobi-matrix-identities} & Unreduced Jacobi recurrence and rational resolvent identities\\
\code{interpolative-decomposition} & Maximal minor of a right singular frame\\
\code{orthogonal-polynomials-measure-def} & Invertible finite Hankel systems\\
\code{hankel-moment-positivity} & Gram quadratic forms and determinant expansion\\
\code{orthogonal-polynomial-zeros} & Jacobi eigenspaces and no-common-root recurrence\\
\code{gauss-christoffel-numbers} & Normalized evaluation vectors\\
\code{quadrature-exactness-characterisation} & Polynomial division by the node polynomial\\
\code{gauss-quadrature-derivative-remainder} & Hermite interpolation and scalar Rolle theory\\
\code{gauss-radau-lobatto} & Modified measures and endpoint factors; signs repaired\\
\code{kernel-polynomials-christoffel} & Exact polynomial division and bidiagonal products\\
\code{formal-orthogonal-polynomials} & Nonsingular bilinear Hankel systems\\
\code{jacobi-fraction-second-kind} & Polynomial difference quotients and resolvents\\
\code{symmetric-measure-even-odd-reduction} & Uniqueness, parity, and the square-map pushforward\\
\code{recurrence-coefficient-bounds} & Rayleigh quotients and Cauchy--Schwarz\\
\code{monic-orthogonal-recurrence} & Compressed multiplication and telescoping identities\\
\code{recurrence-coefficient-continuity} & Continuity of finite inverses and simple real roots\\
\code{gauss-remainder-contour-kernel} & Scalar Cauchy formula after exact resolvent cancellation\\
\code{recurrence-coefficient-moment-sensitivity} & Differentiate finite orthogonality equations\\
\code{christoffel-uvarov-modification} & Finite constrained coefficient systems\\
\code{discrete-measure-jacobi-matrix} & Weighted evaluation matrix and Vandermonde rank\\
\code{rational-gauss-quadrature} & Common-denominator polynomial space\\
\code{companion-matrix} & Cyclic basis and annihilator polynomial\\
\code{psd-hermitian-part} & Kernel inclusion and whitening\\
\code{abs-best-approx-asymptotics} & Parity and finite certificates; Stahl equivalent remains\\
\code{exp-rational-halfline} & Finite Gram proof of 1/3; Halphen equivalent remains\\
\bottomrule
\end{longtable}
\section{Algorithmic entries excluded from this request}
These thirty entries remain in the catalogue, but are outside this note's
mathematical scope. This exclusion is by the advertised task, not by the
area name: a structural theorem in a Krylov area may be included, whereas
an iteration or estimator performance guarantee is excluded.
\begin{longtable}{@{}p{0.57\textwidth}p{0.39\textwidth}@{}}
\toprule
Catalogue identifier & Area\\
\midrule\endhead
\code{tgn-expected-error} & Randomized algorithm guarantees (consumers)\\
\code{etrsvd} & Randomized algorithm guarantees (consumers)\\
\code{etgn} & Randomized algorithm guarantees (consumers)\\
\code{rsvd-tail} & Randomized algorithm guarantees (consumers)\\
\code{nystrom-randomized} & Randomized algorithm guarantees (consumers)\\
\code{block-krylov-lra} & Randomized algorithm guarantees (consumers)\\
\code{hutchpp} & Trace, diagonal and quadratic-form estimation\\
\code{sketch-and-solve-ls} & Randomized solvers (consumers)\\
\code{sketch-and-precondition} & Randomized solvers (consumers)\\
\code{nystrom-pcg} & Randomized solvers (consumers)\\
\code{gmres-def} & Polynomial methods for linear systems\\
\code{gmres-diagonalizable} & Polynomial methods for linear systems\\
\code{gmres-elman} & Polynomial methods for linear systems\\
\code{minres-bound} & Polynomial methods for linear systems\\
\code{fom-def} & Polynomial methods for linear systems\\
\code{chebyshev-iteration} & Polynomial methods for linear systems\\
\code{richardson-iteration} & Polynomial methods for linear systems\\
\code{gmres-bgt} & Field of values and non-normal GMRES theory\\
\code{gmres-any-convergence-curve} & Field of values and non-normal GMRES theory\\
\code{arnoldi-fa-error} & Krylov approximation of f(A)b\\
\code{matvec-query-model} & Krylov bases: grade, Arnoldi and Lanczos relations\\
\code{partial-cholesky-nystrom} & Deterministic structural lemmas\\
\code{qr-algorithm-simultaneous-iteration} & Ritz values and eigenvector angles\\
\code{golub-kahan-bidiagonalization} & Krylov bases: grade, Arnoldi and Lanczos relations\\
\code{cg-recurrence} & Polynomial methods for linear systems\\
\code{lanczos-recurrence-algorithm} & Krylov bases: grade, Arnoldi and Lanczos relations\\
\code{golub-kahan-recurrence} & Krylov bases: grade, Arnoldi and Lanczos relations\\
\code{qr-algorithm-convergence} & Ritz values and eigenvector angles\\
\code{rayleigh-quotient-iteration} & Ritz values and eigenvector angles\\
\code{modified-chebyshev-algorithm} & Orthogonal polynomials, Jacobi matrices and Gauss-type quadrature\\
\bottomrule
\end{longtable}
\section{Source and verification record}
The coverage census and current contracts come from the
\href{https://research.chen.pw/NLAlib/}{NLAlib Atlas} and the
\href{https://github.com/tchen-research/NLAlib/blob/e5725bb94e5bd8e326d58ad219260b0218e74c9e/atlas/atlas.json}{commit-pinned catalogue}.
Reusable public names were checked in the
\href{https://github.com/tchen-research/NLAlib/blob/e5725bb94e5bd8e326d58ad219260b0218e74c9e/atlas/declarations.json}{same snapshot's declaration index}.
The two historical external LRA truncation names are absent from that
index and are explicitly marked unverified in the catalogue; no reuse
claim is made for them.

The proof routes distinguish three obligations: a mathematical identity,
its measure/topology or matrix-coordinate interface, and a final Lean
statement matching the catalogue. A completed first obligation does not
by itself certify either of the others. No catalogue status is changed
by this document.
\end{document}
